% Options for packages loaded elsewhere
% Options for packages loaded elsewhere
\PassOptionsToPackage{unicode}{hyperref}
\PassOptionsToPackage{hyphens}{url}
\PassOptionsToPackage{dvipsnames,svgnames,x11names}{xcolor}
%
\documentclass[
  spanish,
  a4paper,
  oneside]{scrbook}
\usepackage{xcolor}
\usepackage[top=25mm,bottom=25mm,left=30mm,right=20mm]{geometry}
\usepackage{amsmath,amssymb}
\setcounter{secnumdepth}{5}
\usepackage{iftex}
\ifPDFTeX
  \usepackage[T1]{fontenc}
  \usepackage[utf8]{inputenc}
  \usepackage{textcomp} % provide euro and other symbols
\else % if luatex or xetex
  \usepackage{unicode-math} % this also loads fontspec
  \defaultfontfeatures{Scale=MatchLowercase}
  \defaultfontfeatures[\rmfamily]{Ligatures=TeX,Scale=1}
\fi
\usepackage{lmodern}
\ifPDFTeX\else
  % xetex/luatex font selection
\fi
% Use upquote if available, for straight quotes in verbatim environments
\IfFileExists{upquote.sty}{\usepackage{upquote}}{}
\IfFileExists{microtype.sty}{% use microtype if available
  \usepackage[]{microtype}
  \UseMicrotypeSet[protrusion]{basicmath} % disable protrusion for tt fonts
}{}
\makeatletter
\@ifundefined{KOMAClassName}{% if non-KOMA class
  \IfFileExists{parskip.sty}{%
    \usepackage{parskip}
  }{% else
    \setlength{\parindent}{0pt}
    \setlength{\parskip}{6pt plus 2pt minus 1pt}}
}{% if KOMA class
  \KOMAoptions{parskip=half}}
\makeatother
% Make \paragraph and \subparagraph free-standing
\makeatletter
\ifx\paragraph\undefined\else
  \let\oldparagraph\paragraph
  \renewcommand{\paragraph}{
    \@ifstar
      \xxxParagraphStar
      \xxxParagraphNoStar
  }
  \newcommand{\xxxParagraphStar}[1]{\oldparagraph*{#1}\mbox{}}
  \newcommand{\xxxParagraphNoStar}[1]{\oldparagraph{#1}\mbox{}}
\fi
\ifx\subparagraph\undefined\else
  \let\oldsubparagraph\subparagraph
  \renewcommand{\subparagraph}{
    \@ifstar
      \xxxSubParagraphStar
      \xxxSubParagraphNoStar
  }
  \newcommand{\xxxSubParagraphStar}[1]{\oldsubparagraph*{#1}\mbox{}}
  \newcommand{\xxxSubParagraphNoStar}[1]{\oldsubparagraph{#1}\mbox{}}
\fi
\makeatother


\usepackage{longtable,booktabs,array}
\newcounter{none} % for unnumbered tables
\usepackage{calc} % for calculating minipage widths
% Correct order of tables after \paragraph or \subparagraph
\usepackage{etoolbox}
\makeatletter
\patchcmd\longtable{\par}{\if@noskipsec\mbox{}\fi\par}{}{}
\makeatother
% Allow footnotes in longtable head/foot
\IfFileExists{footnotehyper.sty}{\usepackage{footnotehyper}}{\usepackage{footnote}}
\makesavenoteenv{longtable}
\usepackage{graphicx}
\makeatletter
\newsavebox\pandoc@box
\newcommand*\pandocbounded[1]{% scales image to fit in text height/width
  \sbox\pandoc@box{#1}%
  \Gscale@div\@tempa{\textheight}{\dimexpr\ht\pandoc@box+\dp\pandoc@box\relax}%
  \Gscale@div\@tempb{\linewidth}{\wd\pandoc@box}%
  \ifdim\@tempb\p@<\@tempa\p@\let\@tempa\@tempb\fi% select the smaller of both
  \ifdim\@tempa\p@<\p@\scalebox{\@tempa}{\usebox\pandoc@box}%
  \else\usebox{\pandoc@box}%
  \fi%
}
% Set default figure placement to htbp
\def\fps@figure{htbp}
\makeatother


% definitions for citeproc citations
\NewDocumentCommand\citeproctext{}{}
\NewDocumentCommand\citeproc{mm}{%
  \begingroup\def\citeproctext{#2}\cite{#1}\endgroup}
\makeatletter
 % allow citations to break across lines
 \let\@cite@ofmt\@firstofone
 % avoid brackets around text for \cite:
 \def\@biblabel#1{}
 \def\@cite#1#2{{#1\if@tempswa , #2\fi}}
\makeatother
\newlength{\cslhangindent}
\setlength{\cslhangindent}{1.5em}
\newlength{\csllabelwidth}
\setlength{\csllabelwidth}{3em}
\newenvironment{CSLReferences}[2] % #1 hanging-indent, #2 entry-spacing
 {\begin{list}{}{%
  \setlength{\itemindent}{0pt}
  \setlength{\leftmargin}{0pt}
  \setlength{\parsep}{0pt}
  % turn on hanging indent if param 1 is 1
  \ifodd #1
   \setlength{\leftmargin}{\cslhangindent}
   \setlength{\itemindent}{-1\cslhangindent}
  \fi
  % set entry spacing
  \setlength{\itemsep}{#2\baselineskip}}}
 {\end{list}}
\usepackage{calc}
\newcommand{\CSLBlock}[1]{\hfill\break\parbox[t]{\linewidth}{\strut\ignorespaces#1\strut}}
\newcommand{\CSLLeftMargin}[1]{\parbox[t]{\csllabelwidth}{\strut#1\strut}}
\newcommand{\CSLRightInline}[1]{\parbox[t]{\linewidth - \csllabelwidth}{\strut#1\strut}}
\newcommand{\CSLIndent}[1]{\hspace{\cslhangindent}#1}

\ifLuaTeX
\usepackage[bidi=basic,shorthands=off]{babel}
\else
\usepackage[bidi=default,shorthands=off]{babel}
\fi
\ifLuaTeX
  \usepackage{selnolig} % disable illegal ligatures
\fi


\setlength{\emergencystretch}{3em} % prevent overfull lines

\providecommand{\tightlist}{%
  \setlength{\itemsep}{0pt}\setlength{\parskip}{0pt}}



 


\AtBeginDocument{\hypersetup{linkcolor=black}}
% ---------- Paquetes tipograficos y microajustes ----------
\usepackage{microtype}        % mejor interletrado y justificado
\usepackage{csquotes}         % comillas tipograficas
\usepackage{iftex}
\ifPDFTeX\else
  % Seleccion de fuentes solo cuando XeTeX/LuaTeX esta disponible.
  \IfFontExistsTF{TeX Gyre Termes}{
    \setmainfont{TeX Gyre Termes}
  }{
    \IfFontExistsTF{Latin Modern Roman}{\setmainfont{Latin Modern Roman}}{}
  }
  \IfFontExistsTF{TeX Gyre Heros}{
    \setsansfont{TeX Gyre Heros}
  }{
    \IfFontExistsTF{Latin Modern Sans}{\setsansfont{Latin Modern Sans}}{}
  }
  \IfFontExistsTF{TeX Gyre Cursor}{
    \setmonofont{TeX Gyre Cursor}
  }{
    \IfFontExistsTF{Latin Modern Mono}{\setmonofont{Latin Modern Mono}}{}
  }
\fi
\usepackage{enumitem}         % listas compactas
\setlist{itemsep=.2em, topsep=.2em}

% ---------- KOMA-Script: estilo de titulos y espaciado ----------
\KOMAoptions{
  headings=big,
  parskip=half,
  fontsize=12pt,
  appendixprefix=true
}
\setlength{\parindent}{1.5em}
\setlength{\parskip}{0.6em}
\usepackage{indentfirst}
\usepackage{setspace}         % Control del interlineado del documento
\setstretch{1.5}

% ---------- Encabezados y pies (scrlayer-scrpage) ----------
\usepackage[automark,headsepline]{scrlayer-scrpage}
\clearpairofpagestyles
\automark[chapter]{chapter}
\ihead{\pagemark}
\ohead{\itshape\headmark}
\setheadsepline{0.4pt}
\renewcommand*{\chaptermarkformat}{}
\renewcommand*{\chapterpagestyle}{scrheadings}
\pagestyle{scrheadings}
\makeatletter
\let\ps@plain\ps@scrheadings
\makeatother
\cfoot{}

% ---------- Hipervinculos mas sobrios ----------
\usepackage{hyperref}
\hypersetup{
  colorlinks=true,
  linkcolor=black,
  citecolor=[rgb]{0.05,0.2,0.5},
  urlcolor=[rgb]{0.05,0.2,0.5},
  pdfauthor={\@author},
  pdftitle={\@title}
}

% ---------- Leyendas de figuras/tablas ----------
\usepackage[labelfont=bf,textfont=it]{caption}
\captionsetup{
  skip=8pt
}

% ---------- Tabla de contenidos ----------
\KOMAoptions{toc=graduated}
\RedeclareSectionCommand[tocnumwidth=3em]{chapter}
\RedeclareSectionCommand[tocindent=3.25em,tocnumwidth=2.8em]{section}
\RedeclareSectionCommand[tocindent=6.5em,tocnumwidth=2.5em]{subsection}
\makeatletter
\newcommand*{\tocdotfill}{\leavevmode\leaders\hbox to .6em{\hss.\hss}\hfill}
\makeatother
\RedeclareSectionCommand[toclinefill=\tocdotfill]{chapter}
\RedeclareSectionCommand[toclinefill=\tocdotfill]{section}
\RedeclareSectionCommand[toclinefill=\tocdotfill]{subsection}
\renewcommand*{\contentsname}{Tabla de contenido}
\setkomafont{chapterentry}{\normalfont}
\setkomafont{chapterentrypagenumber}{\normalfont}

% ---------- Viudas/Huerfanas y cortes de pagina ----------
\clubpenalty=10000
\widowpenalty=10000
\displaywidowpenalty=10000

% ---------- Entorno abstract para scrbook ----------
\providecommand{\abstractname}{Resumen}
\makeatletter
\@ifundefined{abstract}{
  \newenvironment{abstract}{
    \cleardoublepage
    \thispagestyle{plain}
    \null\vfill
    \begin{center}
      {\bfseries\Large \abstractname\par}
    \end{center}\vspace{1em}
    \begingroup
  }{
    \par\endgroup
    \vfill\null
    \cleardoublepage
  }
}{}
\makeatother

% ---------- Soporte de subtitulo desde YAML ----------
\makeatletter
\providecommand{\subtitle}[1]{\gdef\@subtitle{#1}}
\providecommand{\@subtitle}{}
\makeatother

\makeatletter
\newcommand{\facso@iflist}[4]{%
  \begingroup
    \def\addvspace##1{}%
    \def\numberline##1{##1}%
    \def\facso@target{#2}%
    \newcount\facso@entrycount
    \facso@entrycount=0\relax
    \long\def\contentsline##1##2##3{%
      \def\facso@this{##1}%
      \ifx\facso@this\facso@target
        \advance\facso@entrycount\@ne
      \fi
    }%
    \InputIfFileExists{\jobname.#1}{}{}%
  \endgroup
  \ifnum\facso@entrycount>0\relax
    #3%
  \else
    #4%
  \fi}
\renewcommand{\listoffigures}{%
  \facso@iflist{lof}{figure}{%
    \cleardoublepage
    \phantomsection
    \addchap*{Listado de Figuras}%
    \thispagestyle{scrheadings}%
    \markboth{Listado de Figuras}{Listado de Figuras}%
    \addcontentsline{toc}{chapter}{Listado de Figuras}%
    \@starttoc{lof}%
    \cleardoublepage
    \automark[chapter]{chapter}%
  }{}}
\renewcommand{\listoftables}{%
  \facso@iflist{lot}{table}{%
    \cleardoublepage
    \phantomsection
    \addchap*{Listado de Tablas}%
    \thispagestyle{scrheadings}%
    \markboth{Listado de Tablas}{Listado de Tablas}%
    \addcontentsline{toc}{chapter}{Listado de Tablas}%
    \@starttoc{lot}%
    \cleardoublepage
    \automark[chapter]{chapter}%
  }{}}
\makeatother

% --- Desactivar portada y abstract automaticos de Pandoc (PDF) ---
\AtBeginDocument{\let\maketitle\relax}
\renewenvironment{abstract}{}{}
\providecommand{\appendixname}{}
\providecommand{\appendixtocname}{}
\providecommand{\appendixpagename}{}
\renewcommand*{\appendixname}{Anexo}
\renewcommand*{\appendixtocname}{Anexos}
\renewcommand*{\appendixpagename}{Anexos}
\usepackage{booktabs}
\usepackage{longtable}
\usepackage{array}
\usepackage{multirow}
\usepackage{wrapfig}
\usepackage{float}
\usepackage{colortbl}
\usepackage{pdflscape}
\usepackage{tabularx}
\usepackage{xltabular}
\usepackage{threeparttable}
\usepackage{threeparttablex}
\usepackage[normalem]{ulem}
\usepackage{makecell}
\usepackage{xcolor}
\makeatletter
\@ifpackageloaded{bookmark}{}{\usepackage{bookmark}}
\makeatother
\makeatletter
\@ifpackageloaded{caption}{}{\usepackage{caption}}
\AtBeginDocument{%
\ifdefined\contentsname
  \renewcommand*\contentsname{Tabla de contenidos}
\else
  \newcommand\contentsname{Tabla de contenidos}
\fi
\ifdefined\listfigurename
  \renewcommand*\listfigurename{Listado de Figuras}
\else
  \newcommand\listfigurename{Listado de Figuras}
\fi
\ifdefined\listtablename
  \renewcommand*\listtablename{Listado de Tablas}
\else
  \newcommand\listtablename{Listado de Tablas}
\fi
\ifdefined\figurename
  \renewcommand*\figurename{Lista de figuras}
\else
  \newcommand\figurename{Lista de figuras}
\fi
\ifdefined\tablename
  \renewcommand*\tablename{Lista de tablas}
\else
  \newcommand\tablename{Lista de tablas}
\fi
}
\@ifpackageloaded{float}{}{\usepackage{float}}
\floatstyle{ruled}
\@ifundefined{c@chapter}{\newfloat{codelisting}{h}{lop}}{\newfloat{codelisting}{h}{lop}[chapter]}
\floatname{codelisting}{Listado}
\newcommand*\listoflistings{\listof{codelisting}{Listado de Listados}}
\makeatother
\makeatletter
\makeatother
\makeatletter
\@ifpackageloaded{caption}{}{\usepackage{caption}}
\@ifpackageloaded{subcaption}{}{\usepackage{subcaption}}
\makeatother
\usepackage{bookmark}
\IfFileExists{xurl.sty}{\usepackage{xurl}}{} % add URL line breaks if available
\urlstyle{same}
\hypersetup{
  pdftitle={En un país que envejece: Análisis de la utilización de servicios de salud en Chile},
  pdfauthor={Elisa Calderón; Amalia Tsukame},
  pdflang={es},
  colorlinks=true,
  linkcolor={Maroon},
  filecolor={Maroon},
  citecolor={Blue},
  urlcolor={Blue},
  pdfcreator={LaTeX via pandoc}}


\title{En un país que envejece: Análisis de la utilización de servicios
de salud en Chile}
\usepackage{etoolbox}
\makeatletter
\providecommand{\subtitle}[1]{% add subtitle to \maketitle
  \apptocmd{\@title}{\par {\large #1 \par}}{}{}
}
\makeatother
\subtitle{Informe de avance 2}
\author{Elisa Calderón \and Amalia Tsukame}
\date{30 de octubre de 2026}
\begin{document}
\frontmatter
\maketitle

% --- title-pdf.tex personalizado para portada formal ---
\makeatletter
\providecommand{\subtitle}[1]{\gdef\@subtitle{#1}}
\providecommand{\@subtitle}{}
\providecommand{\frontmattercontext}{}
\providecommand{\advisorname}{}
\providecommand{\advisorlabel}{Profesor guia:}
\providecommand{\frontmatterlocation}{}
\InputIfFileExists{includes/cover-config.tex}{}{}

\newcommand{\PrintTitle}{%
  {\sffamily\bfseries\fontsize{24pt}{28pt}\selectfont \@title\par}%
}
\newcommand{\PrintSubtitle}{%
  \begingroup
  \edef\temp{\detokenize{\@subtitle}}%
  \ifx\temp\empty\relax
    % sin subtitulo
  \else
    {\sffamily\bfseries\large \@subtitle\par}%
  \fi
  \endgroup
}
\newcommand{\PrintAuthor}{%
  {\Large\bfseries Elisa Calder\'on \\[2pt] Amalia Tsukame\par}%
  \vspace{4mm}
  {\small
   Departamento de Sociolog\'ia, Universidad de Chile\\
   ORCID: 0009-0001-0171-7199 \quad\textbar\quad 0009-0004-3578-1289\par}%
}
\newcommand{\PrintDate}{%
  {\small \@date\par}%
}
\makeatother

\begin{titlepage}
\thispagestyle{empty}
\begin{center}
\vspace*{10mm}

% Logo o imagen institucional
\includegraphics[width=0.35\textwidth]{assets/cover.png}\par
\vspace{12mm}

% Titulo y subtitulo
\PrintTitle
\vspace{6mm}
\PrintSubtitle

\vspace{22mm}
\begingroup
\edef\temp{\detokenize{\frontmattercontext}}%
\ifx\temp\empty\relax
  % sin contexto adicional
\else
  {\normalsize \frontmattercontext\par}
\fi
\endgroup

\vspace{18mm}
\PrintAuthor

\vspace{16mm}
\begingroup
\edef\temp{\detokenize{\advisorname}}%
\ifx\temp\empty\relax
  % sin tutor
\else
  \rule{0.45\textwidth}{0.4pt}\par
  {\small \advisorlabel\ \advisorname\par}
\fi
\endgroup

\vfill
\begingroup
\edef\temp{\detokenize{\frontmatterlocation}}%
\ifx\temp\empty\relax
  % sin ubicacion
\else
  {\small \frontmatterlocation\par}
\fi
\endgroup
\PrintDate

\end{center}
\end{titlepage}

% Preliminares (numeros romanos) y estilo simple
\frontmatter
\pagestyle{scrheadings}

\setcounter{tocdepth}{2} % (o 1/3 segun prefieras)
\tableofcontents
\listoffigures
\listoftables


\mainmatter
\bookmarksetup{startatroot}

\chapter*{Presentación}\label{presentaciuxf3n}
\addcontentsline{toc}{chapter}{Presentación}

\markboth{Presentación}{Presentación}

La presente investigación se enmarca en nuestro seminario para optar al
grado de Licenciatura en Sociología de la Universidad de Chile. El
trabajo fue desarrollado en el Departamento de Sociología y cuenta con
la supervisión del profesor guía, Dr.~Juan Carlos Castillo Valenzuela.

\noindent\small\textbf{Palabras clave:} Utilización de servicios de
salud; Envejecimiento poblacional; Desigualdades en salud; Modelo
conductual de Andersen; Análisis multivariado.

\par

\bookmarksetup{startatroot}

\chapter{Introducción}\label{introducciuxf3n}

A nivel global, las dinámicas demográficas han experimentado una
tendencia hacia el envejecimiento. En este contexto, Chile se ha
convertido en uno de los países que ha envejecido más rápido en
comparación con otras regiones del mundo
(\citeproc{ref-unitednationsdepartmentofeconomicandsocialaffairspopulationdivision_world_2022}{United
Nations, 2022}), y atraviesa uno de los procesos de envejecimiento más
avanzados de América Latina
(\citeproc{ref-huenchuan_envejecimiento_2018}{Huenchuan, 2018}). Este
cambio se evidencia en los resultados del Censo 2024, donde la población
de 65 años y más representa el 14\% de la población nacional, frente al
6,6\% registrado en 1992, y el índice de envejecimiento alcanzó 79
personas de 65 años y más por cada 100 personas menores de 15 años
(\citeproc{ref-institutonacionaldeestadisticas_primeros_2025}{Instituto
Nacional de Estadísticas, 2025b}). Este fenómeno también se manifiesta
en la propia población mayor: las personas de 80 años y más, que en 1992
representaban el 11\% de la población adulta mayor, alcanzarían un 28\%
en el año 2050 según las proyecciones (INE,
\citeproc{ref-institutonacionaldeestadisticas_transicion_2022}{2022}).

Esta transformación de la estructura por edades plantea desafíos para la
organización de los sistemas e instituciones públicas y, en particular,
para la planificación de los servicios de salud
(\citeproc{ref-huenchuan_proteccion_2011}{Huenchuan, 2011}), ya que el
aumento de la longevidad y la mayor presencia de personas en edades
avanzadas se producen en un contexto epidemiológico en el que han
adquirido relevancia las enfermedades crónicas no transmisibles, la
multimorbilidad y las limitaciones funcionales
(\citeproc{ref-albala_envejecimiento_2020}{Albala, 2020}). Según la
Encuesta Nacional de Calidad de Vida y Salud (ENCAVI) 2023-2024, el
40,1\% de la población de 15 años y más declara dos o más enfermedades
diagnosticadas, y la declaración de enfermedades es mayor en las edades
avanzadas
(\citeproc{ref-ministeriodesaluddepartamentodeepidemiologia_encuesta_2025}{Ministerio
de Salud, Departamento de Epidemiología, 2025}), lo que proyecta un
aumento en la demanda de los servicios de salud
(\citeproc{ref-huenchuan_proteccion_2011}{Huenchuan, 2011}). No
obstante, es importante reconocer que las necesidades de salud no se
distribuyen de manera homogénea en la población ya que quienes ocupan
mejores posiciones sociales tienen menor probabilidad de enfermar y,
cuando enferman, cuentan con más recursos para enfrentarlo
(\citeproc{ref-moscosocuaresma_determinantes_2017}{Moscoso Cuaresma
et~al., 2017}).

Desde una perspectiva sociológica, esta heterogeneidad puede
comprenderse a partir de la teoría de la ventaja y desventaja acumulada,
según la cual las desigualdades en ingresos, salud y estatus tienden a
divergir con el paso del tiempo
(\citeproc{ref-dannefer_cumulative_2003}{Dannefer, 2003};
\citeproc{ref-ferraro_aging_2009}{Ferraro \& Shippee, 2009}). Desde este
enfoque, personas de la misma edad pueden llegar a la vejez con
distintas necesidades de salud y con posibilidades desiguales de acceder
a la atención. Además, la magnitud de este proceso no es la misma en
todas las sociedades, sino que depende de sus instituciones
(\citeproc{ref-dannefer_cumulative_2003}{Dannefer, 2003}).

En Chile, el sistema de salud está segmentado entre un seguro público,
el Fondo Nacional de Salud (Fonasa), y aseguradoras privadas, las
Isapres, generando una distribución desigual de los recursos y
diferencias en el acceso y calidad de la atención
(\citeproc{ref-bassdelcampo_universalizacion_2025}{Bass del Campo,
2025}). Esto se acompaña de desigualdades territoriales en donde los
médicos especialistas y la infraestructura de mayor complejidad se
concentran en la Región Metropolitana y en los principales centros
urbanos (\citeproc{ref-clinicasdechile_dimensionamiento_2025}{Clínicas
de Chile, 2025};
\citeproc{ref-superintendenciadesalud_caracterizacion_2023}{Superintendencia
de Salud, 2023}), mientras que en las zonas rurales la atención se basa
principalmente en postas de salud que presentan importantes limitaciones
estructurales y de atención (MINSAL y Organización Panamericana de la
Salud {[}OPS{]},
\citeproc{ref-ministeriodesalud_levantamiento_2025}{2025}).

El argumento central de esta investigación es que, en un contexto de
envejecimiento poblacional, la mayor necesidad de atención no se traduce
automáticamente en utilización de servicios de salud: esta también se
asocia con los recursos económicos, de cobertura y las condiciones
territoriales de las personas, y que pueden variar entre grupos de edad.
Frente a este escenario, interesa conocer cómo la utilización de los
servicios de salud se distribuye según los recursos económicos, la
cobertura y el territorio, y si esa distribución cambia a lo largo de la
vida adulta. Para esto, se utilizará el Modelo Conductual de Utilización
de Servicios de Salud de Andersen
(\citeproc{ref-andersen_revisiting_1995}{1995},
\citeproc{ref-andersen_national_2008}{2008}), que ofrece un marco para
examinar estas relaciones al distinguir entre factores predisponentes,
de necesidad y facilitadores de la utilización.

Para abordar empíricamente esta problemática, la presente investigación
utilizará
\href{https://observatorio.ministeriodesarrollosocial.gob.cl/encuesta-casen-2024}{Encuesta
de Caracterización Socioeconómica Nacional (CASEN) 2024} y se analizará
la utilización declarada de tres tipos de servicio en la población de 18
años y más residente en viviendas particulares en Chile: medicina
general, atención de especialidad y hospitalización o intervención
quirúrgica. Se estimarán modelos de regresión logística binaria
anidados, incorporando progresivamente los distintos grupos de factores
de Andersen para examinar su asociación con la utilización de los
servicios.

\section{Objeto de estudio}\label{objeto-de-estudio}

El objeto de estudio corresponde a la utilización declarada de servicios
de salud (atención de medicina general, de especialidad y
hospitalización o intervención quirúrgica) en la población adulta
residente en viviendas particulares en Chile, y su diferenciación según
la edad, el ingreso, la previsión de salud y el territorio de
residencia.

\section{Pregunta de investigación}\label{pregunta-de-investigaciuxf3n}

A partir de lo anterior, se formula la siguiente pregunta de
investigación:

¿Qué factores se asocian con la utilización de servicios de salud en las
distintas edades de la población adulta en Chile?

\section{Objetivo general}\label{objetivo-general}

Analizar qué factores se asocian con la utilización de servicios de
salud en la población adulta de Chile y si esas asociaciones cambian
entre grupos de edad.

\section{Objetivos específicos}\label{objetivos-especuxedficos}

\begin{enumerate}
\def\labelenumi{\arabic{enumi}.}
\tightlist
\item
  Describir los patrones de utilización de medicina general, atención de
  especialidad y hospitalización o intervención quirúrgica según grupo
  de edad en la población de 18 años y más en Chile.
\item
  Estimar la asociación del ingreso y la previsión de salud con la
  utilización de medicina general, atención de especialidad y
  hospitalización o intervención quirúrgica, considerando ambas
  variables conjuntamente.
\item
  Analizar la asociación del territorio de residencia con la utilización
  de medicina general, atención de especialidad y hospitalización o
  intervención quirúrgica.
\item
  Explorar si las asociaciones del ingreso, la previsión y el territorio
  con la utilización de medicina general, atención de especialidad y
  hospitalización o intervención quirúrgica varían entre grupos de edad.
\item
  Explorar si los patrones de utilización y sus asociaciones difieren
  entre hombres y mujeres.
\end{enumerate}

\bookmarksetup{startatroot}

\chapter{Antecedentes conceptuales}\label{antecedentes-conceptuales}

\section{Envejecimiento: transición demográfica y
epidemiológica}\label{envejecimiento-transiciuxf3n-demogruxe1fica-y-epidemioluxf3gica}

El envejecimiento constituye una de las transformaciones sociales más
relevantes del siglo XXI
(\citeproc{ref-unitednationsdepartmentofeconomicandsocialaffairspopulationdivision_world_2020}{United
Nations, 2020}) y se configura como un fenómeno multidimensional, que
puede abordarse desde perspectivas biológicas, sociales, económicas y
culturales.

De manera específica, el envejecimiento admite dos niveles de análisis,
relacionados pero conceptualmente distintos. En el plano individual, es
un proceso gradual, que ocurre a lo largo de toda la vida y es
irreversible (\citeproc{ref-chackiel_envejecimiento_2000}{Chackiel,
2000}). Por otro lado, en el plano demográfico, la Comisión Económica
para América Latina y el Caribe
(\citeproc{ref-comisioneconomicaparaamericalatinayelcaribe_acerca_}{Comisión
Económica para América Latina y el Caribe, s.~f.}) lo define como el
aumento progresivo de la proporción de las personas de 60 años y más con
respecto a la población total, lo que resulta en una progresiva
alteración de la estructura por edades. En Chile, la Ley N° 19.828
define como adulta mayor a quien ha cumplido 60 años y adulto mayor de
la cuarta edad a quien ha cumplido 80 años (\citeproc{ref-_ley_2002}{Ley
N° 19.828, 2002}).

Esta transformación de la estructura etaria se vincula con la transición
demográfica, un proceso histórico desarrollado inicialmente en Europa
para explicar los cambios demográficos, y que describe el paso de
niveles elevados a niveles bajos de natalidad y mortalidad (INE,
\citeproc{ref-institutonacionaldeestadisticas_transicion_2022}{2022}).
Para explicar este proceso, Chesnais (1990, como se citó en
\citeproc{ref-chackiel_envejecimiento_2000}{Chackiel, 2000}) distingue
distintas etapas demográficas: el descenso inicial de la mortalidad, con
una fecundidad todavía elevada, rejuvenece la población; luego, la caída
de la fecundidad reduce el peso relativo de la población infantil; y
cuando la fecundidad alcanza niveles muy bajos y las ganancias en
sobrevida se concentran en las edades avanzadas, el envejecimiento se
intensifica. En la misma línea, la CEPAL clasifica a los países según el
nivel de su fecundidad y la proporción de personas mayores en cinco
etapas de envejecimiento: incipiente, moderado, moderadamente avanzado,
avanzado y muy avanzado
(\citeproc{ref-huenchuan_envejecimiento_2018}{Huenchuan, 2018};
\citeproc{ref-rojas_envejecimiento_}{Hugo Rojas et~al., 2022}).

Aunque la transición demográfica ofrece un marco general, su duración y
su velocidad varían entre regiones. Mientras que en Europa tomó cerca de
dos siglos alcanzar las etapas avanzadas, algunos países de América
Latina y el Caribe las recorrieron en menos de medio siglo (INE,
\citeproc{ref-institutonacionaldeestadisticas_transicion_2022}{2022};
\citeproc{ref-rojas_envejecimiento_}{Hugo Rojas et~al., 2022}), a una
velocidad heterogénea (CEPAL,
\citeproc{ref-comisioneconomicaparaamericalatinayelcaribe_impactos_2026}{2026})
que ha superado la capacidad institucional de las sociedades para
generar condiciones de desarrollo sostenible y protección social
(\citeproc{ref-turra_transicion_2021}{Turra \& Fernandes, 2021}). De
manera específica, en Chile, la fecundidad se encuentra muy por debajo
del nivel de reemplazo generacional (2,1) con 1,03 nacidos vivos
promedio por mujer para el año 2024 (INE,
\citeproc{ref-institutonacionaldeestadisticas_cifras_2025}{2025a}), lo
que ha reducido la proporción de personas de 14 años o menos desde un
29,4\% en 1992 a un 17,7\% en 2024, mientras que la esperanza de vida
supera los 81 años (INE,
\citeproc{ref-institutonacionaldeestadisticas_cifras_2025}{2025a}).
Estos procesos, en su conjunto, ubican al país en la clasificación de
envejecimiento avanzado
(\citeproc{ref-huenchuan_envejecimiento_2018}{Huenchuan, 2018}). Según
las proyecciones del INE, este proceso se profundizará y, para el año
2045, la población de 65 años y más triplicaría a la menor de 15 años
(INE,
\citeproc{ref-institutonacionaldeestadisticas_estimaciones_2026}{2026}).

En paralelo, junto con el proceso de envejecimiento se ha producido en
el país una transición epidemiológica, en la cual se ha experimentado un
aumento de la morbimortalidad por enfermedades no transmisibles, tanto
crónicas como agudas, siendo las primeras causas de muerte para el año
2023 las enfermedades del sistema circulatorio (31.149), neoplasias
(30.796), enfermedades respiratorias (13.728) y causas externas (8.279)
(MINSAL,
\citeproc{ref-ministeriodesaluddepartamentodeepidemiologia_informe_2023}{2023};
\citeproc{ref-departamentodeestadisticaseinformaciondesalud_defunciones_2026}{Departamento
de Estadísticas e Información de Salud, 2026}). Albala
(\citeproc{ref-albala_envejecimiento_2020}{2020}) señala también que el
envejecimiento produce una acumulación de enfermedades crónicas en
mayores de 60 años, constituyéndose como un factor que influye en la
utilización de los servicios de salud
(\citeproc{ref-travassos_uma_2004}{Travassos \& Martins, 2004}).

\section{Curso de vida y desigualdad
acumulada}\label{curso-de-vida-y-desigualdad-acumulada}

Para analizar las diferencias entre personas de distintas edades,
resulta útil el enfoque del curso de vida, que desde finales de 1990 ha
adquirido relevancia en la investigación sobre envejecimiento, al
permitir examinar cómo los eventos históricos y los cambios económicos,
demográficos y sociales configuran tanto las vidas individuales como las
de las cohortes o generaciones
(\citeproc{ref-blanco_enfoque_2011}{Blanco, 2011};
\citeproc{ref-robledomarin_teorias_2021}{Robledo Marín \& Orejuela
Gómez, 2021}). Su concepto central es el de trayectoria: un recorrido a
lo largo de la vida que se despliega en ámbitos interdependientes, como
el trabajo, la educación o la familia, y que está marcado por
transiciones y puntos de inflexión
(\citeproc{ref-blanco_enfoque_2011}{Blanco, 2011}). Desde esta
perspectiva, la vejez no es una etapa aislada, sino el resultado de
trayectorias previas, lo que permite preguntarse por qué personas de una
misma edad llegan a ella en situaciones muy distintas.

La teoría de la ventaja y desventaja acumulada (CAD) ofrece esa
explicación. Dannefer (\citeproc{ref-dannefer_cumulative_2003}{2003}) la
define como una ``tendencia sistémica a la divergencia interindividual''
en características como el dinero, la salud o el estatus a medida que
pasa el tiempo (p.~S327). Su punto central es que esta divergencia no es
una simple proyección de las posiciones iniciales, sino el resultado de
la interacción de múltiples fuerzas, y que es una propiedad de
colectividades, como las cohortes, más que de individuos aislados
(p.~S327). Por eso, las diferencias asociadas a la edad deben leerse
como el producto de procesos que se acumulan a lo largo del curso de la
vida.

Ferraro \& Shippee (\citeproc{ref-ferraro_aging_2009}{2009}) desarrollan
esta idea en la teoría de la desigualdad acumulada (CI). Su primer
planteamiento es que los sistemas sociales generan desigualdad y que
esta no es principalmente resultado de decisiones individuales, sino de
arreglos estructurales. Su segundo planteamiento es que la desventaja
aumenta la exposición a riesgos, mientras que la ventaja aumenta la
exposición a oportunidades, de modo que ambas no son simples opuestos en
una misma escala. Si bien muchos resumen la CAD con la máxima de que
``la ventaja se acumula'', los autores insisten en que lo relevante es
el proceso por el cual esto ocurre y no solo la brecha final. A su vez,
las trayectorias no están determinadas: los recursos disponibles y la
acción de las personas pueden atenuar desventajas previas.

Dannefer (\citeproc{ref-dannefer_cumulative_2003}{2003}) sostiene que el
envejecimiento de cada cohorte implica la reproducción de desigualdades
en salud, con prácticas de salud diferenciadas y acceso estratificado a
la atención (p.~S331). Señala también que dimensiones institucionales,
independientes de las características individuales de las personas,
inciden en factores como los ingresos, la estabilidad laboral y el
acceso a la salud (p.~S332). A esto se suma que la intensidad de la CAD
varía entre sociedades según sus políticas económicas y de bienestar
(p.~S328). Esta idea se puede visualizar para el caso chileno: un
sistema de salud segmentado entre Fonasa e Isapres funciona como una
estructura que distribuye de manera diferenciada las oportunidades de
acceso y utilización
(\citeproc{ref-bassdelcampo_universalizacion_2025}{Bass del Campo,
2025}), de modo que las ventajas y desventajas acumuladas previamente se
traducen en formas distintas de relacionarse con el sistema en la vejez.

Es importante precisar que el diseño transversal de esta investigación
no permite observar directamente la acumulación en el tiempo, sino solo
sus huellas en un momento dado. Ferraro \& Shippee
(\citeproc{ref-ferraro_aging_2009}{2009}) advierten que estudiar
únicamente a personas de edades avanzadas implica trabajar con una
población ya ``truncada'' por la mortalidad, lo que puede dar la falsa
impresión de que la desigualdad disminuye en la vejez. Dannefer
(\citeproc{ref-dannefer_cumulative_2003}{2003}) también cuestiona la
inferencia de trayectorias a partir de comparaciones transversales entre
edades (p.~S328-S329). Por eso, ambas aproximaciones se utilizan aquí
como marco interpretativo y no como hipótesis que se pretenda contrastar
de forma directa.

\section{Necesidad, acceso y utilización de servicios de
salud}\label{necesidad-acceso-y-utilizaciuxf3n-de-servicios-de-salud}

Para comprender cómo se vinculan las diferencias por edad y las
condiciones sociales con el sistema sanitario, es necesario distinguir
entre tres conceptos que están relacionados pero que no son
equivalentes: necesidades de salud, acceso a la atención y utilización
de los servicios de salud. La relevancia de distinguirlos radica en que
las necesidades de salud no se traducen necesariamente en una demanda de
atención y, aun cuando las personas requieran asistencia sanitaria, su
utilización puede depender tanto de sus características individuales y
sociales como de las condiciones en que se organiza la oferta de
servicios.

\subsection{Necesidades de salud}\label{necesidades-de-salud}

Asadi-Lari et~al. (\citeproc{ref-asadi-lari_need_2003}{2003}) señalan
que la necesidad de salud es un concepto cuya definición tiene
implicaciones directas en la formulación de políticas y en la
organización de los servicios, ya que una conceptualización restrictiva
puede limitar el reconocimiento de las necesidades de la población y el
alcance de las respuestas destinadas a satisfacerlas. Desde un enfoque
económico tradicional, la necesidad se define a partir de la capacidad
de beneficiarse de una intervención médica, vinculándola con su
efectividad. Sin embargo, este enfoque puede resultar acotado, ya que se
concentra en los beneficios de la atención y puede dejar de lado
dimensiones personales y sociales relevantes para la salud
(\citeproc{ref-asadi-lari_need_2003}{Asadi-Lari et~al., 2003}).

Desde una perspectiva sociológica, Bradshaw (1972, como se citó en
\citeproc{ref-asadi-lari_need_2003}{Asadi-Lari et~al., 2003}) distingue
cuatro tipos de necesidades: las necesidades normativas, las cuales son
identificadas a partir de criterios profesionales; las necesidades
sentidas que corresponden a lo percibido o experimentado por las
personas; las necesidades expresadas, que son aquellas que se verbalizan
o se traducen en un uso de los servicios de salud; y, por último, las
necesidades comparativas identificadas al contrastar las condiciones de
grupos con características similares que viven en contextos diferentes.

En una línea similar, Andersen
(\citeproc{ref-andersen_revisiting_1995}{1995},
\citeproc{ref-andersen_national_2008}{2008}) plantea una distinción
cercana entre necesidad percibida y necesidad evaluada, permitiendo
distinguir entre la existencia de una necesidad, su reconocimiento por
parte de las personas y su manifestación mediante una demanda de
atención. Sin embargo, una necesidad identificada a partir de criterios
profesionales no necesariamente será percibida o expresada por quien la
experimenta, y la ausencia de utilización no permite concluir, por sí
sola, que no exista una necesidad.

\subsection{Acceso a los servicios}\label{acceso-a-los-servicios}

El acceso a los servicios de salud constituye otra dimensión relevante
para comprender la relación entre las necesidades de atención y la
utilización efectiva de los servicios. Su conceptualización también ha
presentado dificultades, tanto por la amplitud del fenómeno como por la
diversidad de criterios utilizados para definirlo y medirlo.

Una primera línea lo entiende como resultado de la interacción entre
factores de oferta y de demanda. Rodriguez Santana et~al.
(\citeproc{ref-rodriguezsantana_need_2023}{2023}) definen el acceso como
la obtención de la atención sanitaria correcta, en el lugar correcto y
en el momento correcto, reconociendo que el acceso no se limita a la
disponibilidad de servicios ni a su utilización efectiva, sino que
también considera si la atención recibida responde apropiadamente a las
necesidades de salud. Una segunda línea lo concibe como una relación de
ajuste. Penchansky \& Thomas
(\citeproc{ref-penchansky_concept_1981}{1981}) conceptualizan el acceso
como el grado de ajuste entre los usuarios potenciales y los recursos
sanitarios, influido por las características de quienes prestan los
servicios y de quienes los requieren. Esta perspectiva permite reconocer
que el acceso no depende exclusivamente de la existencia física de
establecimientos de salud, sino del vínculo entre las condiciones de
vida de la población y las barreras institucionales, económicas o
geográficas de la red asistencial. Una tercera línea, más integradora,
entiende el acceso como un proceso que comprende las oportunidades de
identificar las necesidades de salud, buscar atención, llegar a los
servicios, obtener atención y beneficiarse de ella
(\citeproc{ref-levesque_patientcentred_2013}{Levesque et~al., 2013}). De
esta manera, el acceso puede verse condicionado por barreras que surgen
en distintas etapas, incluso cuando existen servicios disponibles.

Finalmente, Aday \& Andersen (\citeproc{ref-aday_framework_1974}{1974})
señalan que este concepto ha sido más una idea política que operativa y
abordan esta problemática mediante un marco orientado a examinar la
relación entre las características de la población, las condiciones del
sistema sanitario y la utilización de la atención. Desde esta
perspectiva, el acceso no se limita a la existencia de servicios, sino
que comprende las condiciones que influyen en la posibilidad de obtener
atención sanitaria. En las primeras versiones del modelo de Andersen, el
acceso se entendía como el ingreso a los servicios y la atención
posterior; luego se amplió hacia un concepto multidimensional, compuesto
por el acceso potencial y el acceso realizado
(\citeproc{ref-travassos_uma_2004}{Travassos \& Martins, 2004}).

\subsection{Utilización de los servicios de
salud}\label{utilizaciuxf3n-de-los-servicios-de-salud}

La utilización de los servicios de salud también presenta múltiples
definiciones, y los primeros intentos de desarrollar modelos teóricos al
respecto se remontan a las décadas de 1950 y 1960, pasando por el modelo
de creencias en salud, hasta llegar al modelo de Andersen
(\citeproc{ref-travassos_uma_2004}{Travassos \& Martins, 2004}).

Según Travassos \& Martins (\citeproc{ref-travassos_uma_2004}{2004}), el
núcleo del funcionamiento de los sistemas de salud corresponde a la
utilización de los servicios, y comprende las distintas formas en que
las personas usan las prestaciones sanitarias, incluidas las consultas
médicas, las hospitalizaciones y la realización de exámenes preventivos
y diagnósticos.

Desde el modelo de Andersen
(\citeproc{ref-andersen_national_2008}{2008}), la utilización de los
servicios de salud se define como el acceso realizado o efectivo, el
cual está determinado conjuntamente por las características de la
población y por las propiedades del sistema de salud. En esta misma
línea, Levesque et~al.
(\citeproc{ref-levesque_patientcentred_2013}{2013}) señalan que en la
investigación empírica la utilización suele emplearse como el principal
indicador operativo del acceso, dado que el uso efectivo es directamente
observable y medible, a diferencia del acceso potencial.

En una escala poblacional, la utilización de los servicios de salud
también se ha empleado como un indicador indirecto de las necesidades de
atención. Sin embargo, esta aproximación presenta limitaciones, puesto
que el uso observado no refleja necesariamente la totalidad de las
necesidades existentes en una población. De manera concreta, los
registros del uso de los establecimientos también están condicionados
por la disponibilidad de servicios, las decisiones clínicas y las
alternativas de atención existentes, por lo que no pueden interpretarse
como una medida directa de las necesidades de salud
(\citeproc{ref-asadi-lari_need_2003}{Asadi-Lari et~al., 2003}).

Con todo lo anterior, la utilización está moldeada simultáneamente por
la oferta, la demanda y los recursos individuales, por lo que su
análisis requiere controlar por los distintos factores que median entre
tener una necesidad de salud y lograr atención médica efectiva. Es por
esto que no permite captar por sí sola todas las dimensiones del acceso
a la atención y tampoco permite afirmar que las necesidades fueron
realmente cubiertas al momento de utilizar el sistema de salud.

\section{Modelo Conductual de
Andersen}\label{modelo-conductual-de-andersen}

Para analizar los factores asociados a la utilización de servicios
sanitarios, esta investigación adopta como marco teórico el Modelo
Conductual de Utilización de los Servicios de Salud, desarrollado
inicialmente por Ronald Andersen en 1968 y posteriormente revisado y
ampliado (\citeproc{ref-andersen_revisiting_1995}{Andersen, 1995},
\citeproc{ref-andersen_national_2008}{2008}). El modelo surgió con el
propósito de comprender por qué las personas y familias utilizan los
servicios médicos, establecer criterios operativos para medir el acceso
equitativo a la atención y contribuir al diseño de políticas orientadas
a promover dicho acceso (\citeproc{ref-andersen_national_2008}{Andersen,
2008}).

En su formulación clásica, el modelo plantea que el uso de servicios
depende de la predisposición de las personas a utilizarlos, de los
factores que facilitan o dificultan ese uso y de su necesidad de
atención (\citeproc{ref-andersen_national_2008}{Andersen, 2008}). A
partir de ello distingue tres grupos de factores.

En primer lugar, los factores predisponentes, los que corresponden a
características individuales y sociales que anteceden a la utilización y
estructuran la propensión a utilizar servicios, como son las
características demográficas, la estructura social y las creencias
relacionadas con la salud. Por sí solos no determinan el uso, pero
definen la inclinación de las personas a buscar atención.

En segundo lugar, los factores facilitadores, los cuales representan
recursos y condiciones que pueden permitir o dificultar la utilización,
como los recursos económicos, la cobertura de salud y la disponibilidad
de medios para acceder a la atención. Son relevantes porque representan
la dimensión del acceso que puede ser modificada mediante políticas y
por la organización del sistema.

Por último, los factores de necesidad, los cuales representan las
condiciones que hacen necesaria o potencialmente necesaria la
utilización de servicios. Andersen distingue entre la necesidad
percibida por las propias personas y la necesidad evaluada según
criterios profesionales, distinción cercana a la que plantea Bradshaw
entre necesidades sentidas y normativas.

Estos tres grupos no operan al mismo nivel. En el modelo, los factores
predisponentes anteceden a los demás, los facilitadores habilitan o
restringen que la predisposición se concrete en uso, y la necesidad
constituye el determinante más próximo de la utilización
(\citeproc{ref-travassos_uma_2004}{Travassos \& Martins, 2004}). Para
esta investigación, esto implica que la edad, al ser una característica
predisponente, puede relacionarse con la utilización tanto por las
necesidades de salud que acompañan al envejecimiento como por
características individuales y sociales.

A su vez, esta articulación refleja el desarrollo conceptual del modelo.
En tanto, las primeras versiones explicaban la utilización
principalmente a partir de características individuales, las versiones
posteriores incorporan el sistema de salud, el contexto social, la
interacción entre estos factores y la experiencia previa en el uso de
servicios, de modo que la utilización pasa a entenderse como el
resultado conjunto de factores individuales y contextuales
(\citeproc{ref-andersen_national_2008}{Andersen, 2008};
\citeproc{ref-travassos_uma_2004}{Travassos \& Martins, 2004}).

Andersen (\citeproc{ref-andersen_national_2008}{2008}) señala que estas
características se miden a nivel agregado e incluyen aspectos
relacionados con la comunidad, la organización del sistema sanitario y
los prestadores de servicios. De este modo, los factores predisponentes,
facilitadores y de necesidad pueden examinarse tanto en el plano
individual como en el contextual. A nivel poblacional se pueden
considerar, por ejemplo, características sociodemográficas de la
comunidad, disponibilidad de personal e infraestructura sanitaria e
indicadores de morbilidad, mortalidad y discapacidad. Esta ampliación
permite incorporar al análisis las condiciones del entorno en el que las
personas desarrollan sus necesidades y buscan atención, y en particular
el rol de la oferta sanitaria en los territorios.

El modelo se vincula también con la evaluación de la equidad en el
acceso. Andersen (\citeproc{ref-andersen_national_2008}{2008}) refiere
que la equidad es subjetiva, pero tradicionalmente se ha definido el
acceso equitativo como aquel que se produce cuando las variables
demográficas y de necesidad explican la mayor parte de la varianza en la
utilización, mientras que el acceso inequitativo se produce cuando la
estructura social, las creencias sobre la salud y los recursos
facilitadores determinan quién recibe atención médica. En este estudio,
esta distinción se utiliza como criterio de interpretación, en la cual,
una utilización que varía según factores distintos de la necesidad se
entenderá como indicio de inequidad en el acceso, y no como un efecto
causal.

En este sentido, la relevancia del modelo para el estudio del
envejecimiento radica en que históricamente se ha vinculado al análisis
de encuestas nacionales sobre utilización sanitaria, condiciones
socioeconómicas y estado de salud
(\citeproc{ref-andersen_national_2008}{Andersen, 2008}), lo que lo hace
compatible con el estudio de fenómenos poblacionales. En definitiva,
este modelo proporciona una estructura analítica para estudiar la
utilización de consultas generales, de especialidad y hospitalizaciones
no solo como respuesta a una necesidad individual asociada a la edad,
sino como un fenómeno condicionado por los recursos disponibles y el
contexto socio-sanitario en el que envejecen las personas.

\section{El sistema de salud chileno como contexto: organización,
mecanismos de protección y distribución
territorial}\label{el-sistema-de-salud-chileno-como-contexto-organizaciuxf3n-mecanismos-de-protecciuxf3n-y-distribuciuxf3n-territorial}

Las instituciones constituyen mecanismos mediante los cuales se
canalizan las demandas de los grupos vulnerables y se sostiene la
cooperación colectiva
(\citeproc{ref-huenchuan_envejecimiento_2016}{Huenchuan, 2016}). El
sistema de salud es la institución central para responder a las
necesidades sanitarias de la población, pero como se mencionó
anteriormente, la forma en que organizan y distribuyen sus recursos
puede facilitar o restringir el acceso de distintos grupos sociales.

En el caso de Chile, el sistema de salud se caracteriza por ser un
modelo mixto que abarca tanto el financiamiento y aseguramiento como la
provisión de los servicios médicos, integrando al sector público y
privado (\citeproc{ref-superintendenciadesalud_como_}{Superintendencia
de Salud, s.~f.-a}). Esta estructura dual busca garantizar la cobertura
de toda la población mediante diversos mecanismos regulados y
fiscalizados por el Estado, el cual delega la administración del sistema
en distintas entidades (\citeproc{ref-gonzalez_estructura_2019}{González
et~al., 2019, p. 16}). De esta manera, la rectoría recae en el
Ministerio de Salud, cuya función es formular y controlar las políticas
de salud pública a nivel nacional (MINSAL,
\citeproc{ref-ministeriodesalud_funciones_}{s.~f.}). Por su parte, la
Superintendencia de Salud es un organismo autónomo, dotado de
personalidad jurídica y patrimonio propio, encargado de fiscalizar a
Fonasa, a las Isapres y a los prestadores de salud, y está compuesto por
dos intendencias: la de Fondos y Seguros Previsionales de Salud y la de
Prestadores (\citeproc{ref-gonzalez_estructura_2019}{González et~al.,
2019, p. 41}). Sin embargo, tanto la Organización Panamericana de la
Salud (OPS) como el Ministerio de Salud identifican como principales
deficiencias del sistema su fragmentación administrativa y
organizacional, que se expresa en un financiamiento desarticulado y en
una provisión desigual de recursos
(\citeproc{ref-goldstein_sistema_2018}{Goldstein, 2018}).

El financiamiento de este sistema se sustenta en la obligación legal de
todo trabajador, dependiente e independiente, de cotizar un 7\% de su
renta imponible, permitiéndole al usuario elegir dónde destinar estos
fondos (\citeproc{ref-superintendenciadesalud_como_}{Superintendencia de
Salud, s.~f.-a}). El gran pilar estatal es Fonasa, un seguro público
financiado por estas cotizaciones y aportes fiscales
(\citeproc{ref-gonzalez_estructura_2019}{González et~al., 2019}). Sus
afiliados se dividen en cuatro tramos (A, B, C y D) según sus ingresos;
sin embargo, todos ellos cuentan con atención gratuita (Copago Cero) al
atenderse en la red pública bajo la Modalidad de Atención Institucional
(\citeproc{ref-superintendenciadesalud_como_}{Superintendencia de Salud,
s.~f.-a}). Además, los usuarios de los tramos B, C y D pueden optar por
la Modalidad Libre Elección para comprar bonos y atenderse en clínicas
privadas en convenio
(\citeproc{ref-superintendenciadesalud_como_}{Superintendencia de Salud,
s.~f.-a}).

Por otro lado, se encuentran las Isapres, que administran las
cotizaciones obligatorias, pudiendo también captar cotizaciones
adicionales o voluntarias
(\citeproc{ref-gonzalez_estructura_2019}{González et~al., 2019, p. 42}).
Estas entidades ofrecen contratos bajo lógicas de mercado (estrictamente
reguladas) que se financian con el 7\% legal más una cotización
adicional voluntaria si el plan lo requiere
(\citeproc{ref-gonzalez_estructura_2019}{González et~al., 2019, p. 82}).
De forma paralela a estos dos grandes ejes, existe un subsistema cerrado
operado por la Caja de Previsión de la Defensa Nacional (CAPREDENA),
para las Fuerzas Armadas, y por la Dirección de Previsión de Carabineros
(DIPRECA), a la que adscriben Carabineros, Gendarmería y la PDI
(\citeproc{ref-gonzalez_estructura_2019}{González et~al., 2019}).

En cuanto a la provisión de los servicios, la atención médica se
articula a través de dos grandes redes. La red pública es administrada
por el Estado mediante el Sistema Nacional de Servicios de Salud (SNSS)
y sus 29 divisiones territoriales, que disponen de una red asistencial,
constituida por hospitales e institutos de especialidad
(\citeproc{ref-gonzalez_estructura_2019}{González et~al., 2019, p. 45}).
A su vez la atención sanitaria se organiza en tres niveles según la
complejidad y la cobertura de los problemas de salud que resuelven:
primario, secundario y terciario
(\citeproc{ref-vignolo_niveles_2011}{Vignolo et~al., 2011}).

La atención primaria de salud (APS) es la puerta de entrada al sistema,
tiene una complejidad básica y una amplia cobertura, con atención
ambulatoria. Desarrolla un enfoque preventivo, promocional y familiar;
sus establecimientos están mayoritariamente bajo administración
municipal y su alcance es territorial y comunitario
(\citeproc{ref-gattini_sistema_2018}{Gattini, 2018};
\citeproc{ref-olguin_niveles_2025}{Olguín, 2025}). Está constituida por
Centros de Salud Familiar (CESFAM), Centros Comunitarios de Salud
Familiar (CECOSF), Postas Rurales (PR), Servicios de Atención Primaria
de Urgencia (SAPU) y Servicios de Atención de Urgencia de Alta
Resolutividad (SAR)
(\citeproc{ref-serviciodesaludmetropolitanosur_rol_}{Servicio de Salud
Metropolitano Sur, s.~f.}).

El nivel secundario cuenta con una complejidad intermedia y una menor
cobertura, y requiere equipos interdisciplinarios
(\citeproc{ref-olguin_niveles_2025}{Olguín, 2025}). Sus establecimientos
son los Centros de Referencia de Salud (CRS) y los Centros de
Diagnóstico y Tratamiento (CDT) (SSMS,
\citeproc{ref-serviciodesaludmetropolitanosur_rol_}{s.~f.}).

Por último, el nivel terciario presenta una mayor complejidad, pero una
cobertura reducida de la población. Está constituido por los
establecimientos hospitalarios (SSMS,
\citeproc{ref-serviciodesaludmetropolitanosur_rol_}{s.~f.}), que pueden
recibir pacientes de su propia área de red asistencial o tener un
carácter regional o suprarregional, y cuenta con recursos humanos de
alta especialización y con recursos materiales y de infraestructura de
mayor complejidad técnica (\citeproc{ref-gattini_sistema_2018}{Gattini,
2018}).

Como contraparte, la red privada está conformada por clínicas, centros
médicos, laboratorios y consultas particulares que absorben
principalmente a los afiliados de Isapres, así como a los pacientes de
Fonasa que utilizan la libre elección o que son derivados mediante
licitaciones del Estado (\citeproc{ref-equipoeunamed_como_2025}{Equipo
EUNAMed, 2025}).

Más allá de la segmentación entre redes y seguros, el sistema contempla
protecciones que se aplican a todas las personas. Según la
Superintendencia de Salud
(\citeproc{ref-superintendenciadesalud_garantias_}{s.~f.-b}),
independientemente de si la persona está afiliada a Fonasa o a una
Isapre, existen garantías universales transversales, entre las que se
destaca el Plan de Garantías Explícitas en Salud (GES), que asegura por
ley el acceso, la oportunidad de atención, la calidad y la protección
financiera (con topes máximos de pago) para 90 problemas de salud de
alto impacto. Sin embargo, estas garantías cubren un conjunto acotado de
problemas de salud, y la evidencia previa indica que la implementación
del régimen del Plan de Acceso Universal a Garantías Explícitas (AUGE),
luego denominado GES, aumentó la utilización general sin eliminar las
desigualdades socioeconómicas (\citeproc{ref-nunez_equity_2013}{Núñez \&
Chi, 2013}). Asimismo, la Ley de Urgencia establece que, ante una
emergencia de riesgo vital o riesgo de secuela funcional grave, el
paciente debe ser atendido en el centro de salud más cercano, sea
público o privado, sin que se le pueda exigir un cheque en garantía o
pago previo; esto no implica que la atención sea gratuita, ya que los
costos se pagan posteriormente a través del seguro correspondiente. Una
vez estabilizado, el usuario es derivado a la red que corresponda según
su seguro
(\citeproc{ref-bibliotecadelcongresonacionaldechile_ley_}{Biblioteca del
Congreso Nacional de Chile, s.~f.}).

Finalmente, el contexto regulatorio del sistema se ha visto
profundamente modificado tras la entrada en vigencia de la Ley Corta de
Isapres (Ley N° 21.674) en 2024. Esta normativa fue promulgada para dar
cumplimiento a un fallo de la Corte Suprema, que ordenó aplicar una
Tabla Única de Factores (TUF) sin discriminación por sexo o edad, y para
evitar una quiebra sistémica del sector privado. En ese marco,
estableció la devolución gradual de los cobros en exceso realizados
desde el año 2020. Además, impuso un ajuste de los planes al 7\% legal,
reduciendo los excedentes al exigir mayores beneficios cuando el costo
original era inferior a ese mínimo
(\citeproc{ref-goldstein_resena_2024}{Goldstein, 2024}).

\section{Distribución territorial de la
oferta}\label{distribuciuxf3n-territorial-de-la-oferta}

En el marco del análisis de los sistemas de salud es necesario
comprender cómo las características geográficas condicionan el acceso a
la atención médica, y para esto resulta fundamental precisar qué se
entiende por oferta sanitaria y sus implicancias.

Rodriguez Santana et~al.
(\citeproc{ref-rodriguezsantana_need_2023}{2023}) definen operativamente
la oferta como los servicios e intervenciones preventivas y curativas
proporcionadas formalmente por el sistema sanitario, excluyendo de su
alcance el cuidado informal y la asistencia social comunitaria. A partir
de esta delimitación, el análisis de la oferta no se reduce únicamente
al número de establecimientos existentes, sino a la capacidad instalada
del sistema, comprendiendo infraestructura, camas, tecnología y la
dotación formal de recursos humanos médicos y de especialidad para
responder a las necesidades de la población en un contexto geográfico
determinado. Esta capacidad instalada se distribuye físicamente sobre un
territorio, el cual constituye el espacio concreto donde las personas
acceden a los bienes y servicios que requieren para satisfacer sus
necesidades, ejercer sus derechos sociales y desarrollar sus proyectos
de vida
(\citeproc{ref-programadelasnacionesunidasparaeldesarrollo_desigualdad_2018}{Programa
de las Naciones Unidas para el Desarrollo, 2018}).

En este marco, la disponibilidad de profesionales de la salud, y en
particular de especialistas, es una de las principales brechas
regionales percibidas por la población (PNUD,
\citeproc{ref-programadelasnacionesunidasparaeldesarrollo_desigualdad_2018}{2018}).
A nivel nacional, Chile registra 3,3 médicos por 1.000 habitantes en
2023, por debajo del promedio de la OCDE, que es de 3,7
(\citeproc{ref-bancomundial_medicos_}{Banco Mundial, s.~f.-b};
Organisation for Economic Co-operation and Development {[}OECD{]},
\citeproc{ref-organisationforeconomicco-operationanddevelopment_health_2025}{2025}).
De acuerdo con la Superintendencia de Salud
(\citeproc{ref-superintendenciadesalud_caracterizacion_2023}{2023}), el
39,2\% de los médicos tenía en 2022 un vínculo, no necesariamente
exclusivo, con el sector público; es decir, alrededor del 60\% no
mantenía vínculo con él, y se accede a ellos por vías distintas a la red
pública (\citeproc{ref-clinicasdechile_dimensionamiento_2025}{Clínicas
de Chile, 2025}).

Para el año 2024, el 67,3\% de los médicos en Chile tienen alguna
especialidad, del total, 37,6\% pertenece al sector público y el 62,4\%
al sector privado. La tasa de médicos especialistas es de 1,73 por 1.000
habitantes, ocupando el quinto lugar en la OCDE
(\citeproc{ref-clinicasdechile_dimensionamiento_2025}{Clínicas de Chile,
2025}). En relación con su distribución territorial, la Superintendencia
de Salud
(\citeproc{ref-superintendenciadesalud_caracterizacion_2023}{2023})
muestra diferencias marcadas: la Región Metropolitana registra una tasa
de 2,46 médicos especialistas por cada 1.000 habitantes, seguida por
Biobío (1,72) y Valparaíso (1,43), mientras que Tarapacá y Arica y
Parinacota presentan tasas de 0,70 y 0,81, respectivamente, siendo
alrededor de un tercio de la Región Metropolitana. Estas diferencias
muestran que la disponibilidad territorial de especialistas es
heterogénea y está asociada con factores demográficos y con la capacidad
formadora de las instituciones presentes en aquellas regiones
(\citeproc{ref-huaiquilaf-jorquera_medicos_2026}{Huaiquilaf-Jorquera
et~al., 2026}).

En relación con los establecimientos hospitalarios, para el año 2024 el
62,6\% son establecimientos públicos (SNSS y FFAA), 37,4\% son
prestadores privados
(\citeproc{ref-clinicasdechile_dimensionamiento_2025}{Clínicas de Chile,
2025}). La disponibilidad de camas a nivel nacional es de 1,9 por cada
1.000 habitantes para el año 2023
(\citeproc{ref-bancomundial_camas_}{Banco Mundial, s.~f.-a}). En el
sector privado para el año 2024, la oferta se encuentra fuertemente
centralizada ya que el 64,3\% de las camas están en la Región
Metropolitana, el 19,9\% en la zona sur, y el 15,8\% restante se
encuentra en la zona norte
(\citeproc{ref-clinicasdechile_dimensionamiento_2025}{Clínicas de Chile,
2025}). Esta pronunciada concentración geográfica e institucional de la
infraestructura de alta complejidad configura un acceso dispar a la cama
de un hospital debido a su localización
(\citeproc{ref-fuenzalida_analisis_2023}{Fuenzalida et~al., 2023}), lo
que puede restringir la utilización de servicios hospitalarios y de
especialidad.

En las zonas rurales, la oferta sanitaria enfrenta limitaciones
adicionales. Si bien la ruralidad es variada y en Chile coexisten
distintos criterios para definirla (modificando la magnitud de la
población rural), estos territorios comparten el aislamiento, la
dispersión de la población y una mayor incidencia de pobreza
multidimensional, que en 2022 alcanzó un 27,96\% frente a un 15,54\% en
las zonas urbanas, además de un envejecimiento de su población (MINSAL y
OPS, \citeproc{ref-ministeriodesalud_levantamiento_2025}{2025}).

Para responder a esta realidad territorial, la oferta formal del nivel
primario se apoya sobre las Postas de Salud Rural (PSR), presentes en el
78\% de las comunas del país, de las cuales el 63\% de estas comunas
están clasificadas como rurales. Estos establecimientos constituyen el
primer punto de contacto con las comunidades rurales (MINSAL y OPS,
\citeproc{ref-ministeriodesalud_levantamiento_2025}{2025}). Sin embargo,
la capacidad real de esta oferta sanitaria rural presenta limitaciones
estructurales dado que un 44\% de las PSR a nivel nacional posee una
antigüedad igual o superior a 30 años y un 14\% no cuenta con box de
procedimientos. Asimismo, la comunicación de estos establecimientos es
limitada: el 70\% de las PSR no cuenta con teléfono fijo y el 89\%
carece de teléfono satelital (MINSAL y OPS,
\citeproc{ref-ministeriodesalud_levantamiento_2025}{2025}).

Según el modelo de Andersen
(\citeproc{ref-andersen_national_2008}{2008}), estas características
pueden entenderse como factores facilitadores contextuales que
dificultan el uso de servicios. Dado que la oferta sanitaria rural es
limitada y la distancia a los centros urbanos implica barreras de
entrada a la atención, se hace necesario trasladarse para acceder a
otros niveles de atención, lo que supone un gasto de bolsillo (MINSAL y
OPS, \citeproc{ref-ministeriodesalud_levantamiento_2025}{2025}). Estas
desigualdades territoriales adquieren relevancia para abordar la
distribución desigual de los recursos asociado al modelo de
administración actual y fortalecer la capacidad resolutiva del nivel
primario de atención en salud
(\citeproc{ref-goldstein_sistema_2018}{Goldstein, 2018}).

En síntesis, la distribución de la oferta sanitaria en Chile no es
equilibrada: la concentración de especialistas y de camas, junto con las
limitaciones de la red rural, sugiere que el lugar de residencia puede
modificar las posibilidades de utilizar servicios de especialidad y
hospitalarios, más allá de la necesidad de atención de cada persona.
Asimismo, es relevante considerar que la oferta de servicios no es lo
mismo que el acceso ni que la utilización: que existan especialistas o
camas en un territorio no asegura que todas las personas puedan llegar a
ellos, ni que quienes llegan los utilicen. Por eso, en esta
investigación, la oferta se entiende como una condición del contexto que
puede facilitar o restringir la utilización, y no como una medida de
esta.

\bookmarksetup{startatroot}

\chapter{Antecedentes empíricos}\label{antecedentes-empuxedricos}

\section{Antecedentes empíricos
internacionales}\label{antecedentes-empuxedricos-internacionales}

La literatura internacional ha documentado la relación entre el estatus
socioeconómico y el uso de servicios de salud. En una revisión de
carácter comparado, Lantz
(\citeproc{ref-lantz_socioeconomic_2001}{2001}) señala que, si bien la
disponibilidad y la utilización de la atención de salud dependen de
múltiples factores, los indicadores de estatus socioeconómico (ESE),
como el ingreso, el nivel educacional y la condición ocupacional, se
asocian significativamente con la utilización de servicios en distintos
países. Esta relación abarca tanto la cobertura de seguro como el uso de
servicios personales: consultas de atención primaria y de especialidad,
hospitalizaciones, atención dental y oftalmológica, salud mental,
medicamentos con receta y atención domiciliaria
(\citeproc{ref-lantz_socioeconomic_2001}{Lantz, 2001}).

Con algunas excepciones, la revisión concluye que la mayor parte de los
hallazgos empíricos apunta a una asociación positiva: a mayor ESE, mayor
utilización de atención de salud, y menores niveles de enfermedad sin
tratar y de necesidades médicas insatisfechas. Este patrón se ha
observado tanto en hombres como en mujeres, y tanto en población
infantil como adulta (\citeproc{ref-lantz_socioeconomic_2001}{Lantz,
2001}). En esa línea, los estratos económicos y educativos más bajos
tienden a realizar menos consultas, tanto de atención primaria como de
especialidad, a utilizar menos los servicios de salud mental, dental,
oftalmológica y preventiva, a contar con menos frecuencia con una fuente
habitual de atención y a presentar tasas significativamente mayores de
necesidades de salud no cubiertas. En algunos análisis, un ESE más bajo
se asocia con un mayor uso de servicios médicos y hospitalarios,
especialmente por causas evitables; sin embargo, este patrón parece
explicarse principalmente por la mayor prevalencia de mal estado de
salud en las poblaciones desfavorecidas
(\citeproc{ref-lantz_socioeconomic_2001}{Lantz, 2001}).

Respecto de los factores demográficos, la revisión encuentra que las
edades extremas y el sexo femenino utilizan más los servicios de salud,
efecto que está mediado, al menos en parte, por una mayor necesidad de
atención. El tamaño y la estructura familiar también se asocian a la
utilización, pero la dirección del efecto depende del país: por ejemplo,
en sistemas de acceso universal, los hijos de madres solteras consultan
más, mientras que en sistemas con mayores barreras de acceso, este mismo
grupo consulta en menor proporción, lo que evidencia la desigualdad
(\citeproc{ref-mendoza-sassi_utilizacion_2001}{Mendoza-Sassi \& Béria,
2001}).

El tipo de sistema de salud también modula el efecto del nivel
socioeconómico. En países con acceso universal, como Inglaterra o
Canadá, las personas de clase baja presentan mejor acceso a los
servicios, de modo que el sistema opera como un modificador de efecto
que promueve la equidad. En otros contextos, en cambio, los sistemas no
logran reducir las desigualdades
(\citeproc{ref-mendoza-sassi_utilizacion_2001}{Mendoza-Sassi \& Béria,
2001}).

Por último, entre los factores propios de los servicios de salud, la
menor distancia al establecimiento y el mayor número de médicos por
habitante aumentan la utilización. En países con una red de servicios
adecuada estos factores pierden relevancia, aunque la distancia
geográfica puede volver a ser una barrera para grupos específicos, como
las personas mayores o quienes tienen alguna limitación funcional
(\citeproc{ref-mendoza-sassi_utilizacion_2001}{Mendoza-Sassi \& Béria,
2001}).

\section{Antecedentes empíricos
nacionales}\label{antecedentes-empuxedricos-nacionales}

La investigación nacional ha documentado desigualdades en la utilización
de servicios de salud en Chile, incluso considerando las diferencias en
las necesidades sanitarias de la población. Sin embargo, estas
desigualdades no se expresan de igual manera entre los tipos de
utilización.

\subsection{Facilitadores individuales: ingreso y previsión de
salud}\label{facilitadores-individuales-ingreso-y-previsiuxf3n-de-salud}

En relación a las desigualdades socioeconómicas en la utilización de
servicios de salud, Núñez \& Chi
(\citeproc{ref-nunez_equity_2013}{2013}) analizaron la utilización y la
equidad en el acceso mediante datos de las encuestas CASEN de 1992 a
2009 y de la Encuesta de Satisfacción y Gasto de Bolsillo en Salud de
2006. Sus resultados muestran que, tras ajustar por necesidad, el 20\%
de mayores ingresos utilizaba servicios médicos casi el doble que el
20\% de menores ingresos. En particular, el quintil más pobre registraba
una utilización un 7,3\% inferior a la esperada según sus necesidades,
mientras que en el quintil más rico era un 10,8\% superior a lo
esperado.

En la misma línea y específicamente en las consultas de especialidad,
Vásquez et~al. (\citeproc{ref-vasquez_incomerelated_2013}{2013}), con
las encuestas CASEN de 2000, 2003 y 2009, consideraron variables como la
salud autoevaluada, las limitaciones físicas, la edad y el sexo para
distinguir las diferencias de utilización de aquellas asociadas a
necesidades sanitarias. Encontraron que las consultas de especialidad
presentaban una distribución favorable a los grupos de mayores ingresos,
incluso después de ajustar por necesidad, y que la educación y otros
factores socioeconómicos contribuyeron a explicar estas desigualdades
(pp.~102--103). Paraje \& Vásquez
(\citeproc{ref-paraje_health_2012}{2012}), con la CASEN 2003 y 2009,
también observaron que en ambos años las consultas de especialidad se
concentraron en los grupos de mayores ingresos. Por último, Chovar Vera
et~al. (\citeproc{ref-chovarvera_desigualdad_2014}{2014}) encontraron
que las consultas de especialidad favorecieron a los grupos de mayores
ingresos. En función de lo anterior, se plantea que:

\textbf{H1.} A mayor ingreso, mayor probabilidad de consulta de
especialidad.

En cuanto a la previsión de salud, Henríquez \& Velasco
(\citeproc{ref-henriquez_desigualdades_2015}{2015}) con datos de la
CASEN 1992, 2000 y 2013, estimaron que una mayor probabilidad de
afiliación a Isapre se asociaba con una mayor probabilidad de consulta
de especialidad y de hospitalización. Asimismo, Vásquez et~al.
(\citeproc{ref-vasquez_incomerelated_2013}{2013}) señalan que la
afiliación a una Isapre, junto con el nivel educacional, contribuyó a
explicar las desigualdades en la atención de especialidad
(pp.~102--103). Esto es coherente con un sistema segmentado que
distribuye de manera diferenciada las oportunidades de acceso
(\citeproc{ref-bassdelcampo_universalizacion_2025}{Bass del Campo,
2025}). Por lo tanto, a partir de esto se plantea:

\textbf{H2.} Estar afiliado a Isapre se asocia con mayor probabilidad de
consulta de especialidad y de hospitalización o intervención quirúrgica
que estar afiliado a Fonasa.

\subsection{Facilitadores
contextuales}\label{facilitadores-contextuales}

En Chile, la literatura sobre ruralidad se centra en las diferencias de
acceso, distancia y distribución de los establecimientos de salud. Según
los datos de la CASEN 2022, el 36\% de los hogares rurales con personas
mayores se encuentra a menos de 2,5 km de un centro de salud, frente al
88,9\% de los hogares urbanos
(\citeproc{ref-observatoriodelenvejecimientouc-confuturo_envejecer_2023}{Observatorio
del Envejecimiento UC-Confuturo, 2023}). Si bien, estos resultados
señalan personas con edades avanzadas, ilustran diferencias
territoriales en la proximidad de los servicios. A esto se suma la
variabilidad geográfica del país, por ejemplo, en la macrozona norte
llegar a centros de salud primaria puede tomar entre 57 y 224,3 minutos
y para llegar al hospital o clínica más cercana toma entre 92,9 y 484,9
minutos
(\citeproc{ref-observatoriodelenvejecimientouc-confuturo_envejecer_2023}{Observatorio
del Envejecimiento UC-Confuturo, 2023}). En este contexto, la desigual
distribución de la infraestructura hospitalaria afecta la posibilidad de
acceder a atención en salud, especialmente lejos de la zona central del
país (\citeproc{ref-fuenzalida_analisis_2023}{Fuenzalida et~al., 2023}).

Respecto del uso efectivo de los servicios de salud, la evidencia es
menos concluyente. Paraje \& Vásquez
(\citeproc{ref-paraje_health_2012}{2012}) identificaron que, en 2009, el
promedio de consultas de especialidad en zonas rurales fue menos de la
mitad del observado en zonas urbanas (0,107 frente a 0,237), mientras
que la brecha en medicina general fue menor (0,210 frente a 0,265)
(pp.~12--13, tabla 5). En cambio, Henríquez \& Velasco
(\citeproc{ref-henriquez_desigualdades_2015}{2015}) estimaron que la
zona rural no afectaba la probabilidad de consulta general en 1992 y
2000, pero sí en 2013, mientras que en especialidad y urgencia el efecto
existía en los años anteriores y desapareció en 2013. Por su parte, con
datos de la CASEN 2017, el Proyecto NODO
(\citeproc{ref-proyectonodo_comunidades_2021}{2021}) informó que el 91,7
\% de las personas mayores rurales que reportaron un problema de salud
en los últimos tres meses recibió atención médica, frente al 93,3 \% de
las urbanas, diferencia que no resultó estadísticamente significativa.

Por tanto, aunque la evidencia confirma la relevancia de las barreras
territoriales, sus asociaciones con la utilización varían según el
servicio y el período estudiado. Por esto, la diferencia observada en
las consultas de especialidad resulta especialmente relevante ya que la
oferta de especialistas tiende a concentrarse en los centros urbanos.
Sobre esta base se plantea que:

\textbf{H3.} Residir en zona rural o lejos de un centro de salud se
asocia con menor probabilidad de consulta de especialidad.

\subsection{Medicina general y
hospitalización}\label{medicina-general-y-hospitalizaciuxf3n}

A diferencia de lo observado en las consultas de especialidad, los
resultados sobre medicina general y hospitalización son menos
consistentes. Vásquez et~al.
(\citeproc{ref-vasquez_incomerelated_2013}{2013}) observaron que los
resultados variaban según la medida de utilización: a igual necesidad,
los grupos de menores ingresos concentraban más días de hospitalización,
aunque la probabilidad de hospitalizarse mostraba una distribución
ligeramente favorable a los grupos de mayores ingresos (p.~103, tabla
4). Paraje \& Vásquez (\citeproc{ref-paraje_health_2012}{2012})
encontraron que la hospitalización presentó una distribución favorable a
los grupos de menores ingresos, y que la medicina general pasó de una
distribución ligeramente favorable a ellos en 2003 a otra favorable a
los de mayores ingresos en 2009 (p.~13, tabla 6). Por su parte,
Henríquez \& Velasco (\citeproc{ref-henriquez_desigualdades_2015}{2015})
no encontraron una diferencia significativa en la hospitalización entre
los mayores de 66 años y el tramo de 15 a 25 años (p.~17). Esto es
coherente con la revisión internacional, donde el mayor uso hospitalario
de los grupos desfavorecidos se explica principalmente por su peor
estado de salud.Como la evidencia no permite anticipar una dirección
clara, se exploran:

\textbf{P1.} Se explora si la edad y el ingreso se asocian con la
hospitalización o intervención quirúrgica.

\textbf{P2.} Se explora si el territorio se asocia con la medicina
general y la hospitalización de forma distinta a como lo hace con la
especialidad.

\subsection{Edad y sexo}\label{edad-y-sexo}

En relación con la edad, Chovar Vera et~al.
(\citeproc{ref-chovarvera_desigualdad_2014}{2014}) analizaron la
utilización de servicios médicos con las encuestas CASEN de 2000, 2003,
2009 y 2011, distinguiendo cuatro grupos etarios, sin embargo, en
medicina general y días de hospitalización no se identificó un patrón
consistente entre grupos etarios. Por su parte, Henríquez \& Velasco
(\citeproc{ref-henriquez_desigualdades_2015}{2015}) estimaron que, en
los mayores de 66 años, la probabilidad de consulta general y de
especialidad fue mayor que en el tramo de 15 a 25 años.

Estos antecedentes muestran que la utilización y su desigualdad varían
con la edad, pero no permiten saber si las asociaciones del ingreso, la
previsión y el territorio cambian a lo largo de la vida adulta. Desde la
ventaja y desventaja acumulada cabría esperar que las diferencias sean
mayores en edades avanzadas, porque han tenido más tiempo para
acumularse. Sin embargo, estudiar a la población adulta mayor puede
subestimar la desigualdad ya que la mortalidad actúa como filtro
selectivo en grupos más vulnerables
(\citeproc{ref-ferraro_aging_2009}{Ferraro \& Shippee, 2009}). La
evidencia internacional sugiere, además, que la distancia a los
servicios puede ser una barrera particular para las personas mayores o
con limitación funcional
(\citeproc{ref-mendoza-sassi_utilizacion_2001}{Mendoza-Sassi \& Béria,
2001}). Como el diseño transversal no permite distinguir entre estas
lecturas, se exploran:

\textbf{P3.} Se explora si la asociación entre el ingreso y la
utilización de cada tipo de atención cambia según el grupo de edad.

\textbf{P4.} Se explora si la asociación entre el territorio y la
utilización de cada tipo de atención cambia según el grupo de edad.

\textbf{P5.} Se explora si la asociación entre la previsión y la
utilización de cada tipo de atención cambia según el grupo de edad.

Finalmente, se explora si el sexo modifica estos patrones:

\textbf{P6.} Se explora si los patrones de utilización y su asociación
con el ingreso, la previsión y el territorio difieren entre hombres y
mujeres.

\bookmarksetup{startatroot}

\chapter{Metodología}\label{metodologuxeda}

\section{Datos}\label{datos}

La investigación tendrá un enfoque cuantitativo, observacional y
transversal. Se utilizará como fuente principal de datos la
\href{https://observatorio.ministeriodesarrollosocial.gob.cl/encuesta-casen-2024}{Encuesta
de Caracterización Socioeconómica Nacional (CASEN) 2024}, desarrollada
por el Ministerio de Desarrollo Social y Familia. La encuesta recopila
información sobre aspectos demográficos, módulos de salud, educación,
vivienda, trabajo e ingresos, que permite conocer periódicamente la
situación de los hogares y de la población, especialmente de aquella en
situación de pobreza y de aquellos grupos definidos como prioritarios
por la política social, además de evaluar el impacto de la política
social
(\citeproc{ref-ministeriodedesarrollosocialyfamilia_informe_2024}{Ministerio
de Desarrollo Social y Familia, 2024}). La base de microdatos y su
documentación metodológica son de acceso público a través del
Observatorio Social del MDSF.

La población objetivo de la encuesta corresponde a las personas que
residen en viviendas particulares ocupadas del territorio nacional.
Quedan excluidas las personas que residen en viviendas colectivas, como
establecimientos de larga estadía para personas mayores u hospitales, y
las áreas clasificadas como especiales por el INE en el Marco Muestral
de Viviendas (MMV) 2021, entre ellas diez comunas completas
(\citeproc{ref-ministeriodedesarrollosocialyfamilia_informe_2024}{Ministerio
de Desarrollo Social y Familia, 2024, p. 29}).

El diseño muestral es probabilístico, estratificado y bietápico. Los
estratos se definieron por la combinación de comuna, área y nivel
socioeconómico de las Unidades Primarias de Muestreo (UPM); en cada
estrato, la selección de UPM se realizó de forma independiente, y la
unidad final de selección es la vivienda
(\citeproc{ref-ministeriodedesarrollosocialyfamilia_informe_2024}{Ministerio
de Desarrollo Social y Familia, 2024}). La recolección de datos se
realizó entre el 1 de noviembre de 2024 y el 2 de febrero de 2025,
mediante entrevistas presenciales asistidas por dispositivos
electrónicos a una persona informante integrante del hogar de 18 años o
más, que cuenta con la información necesaria para completar la encuesta,
quien responde la totalidad de los módulos del cuestionario en
representación de todos los integrantes del hogar, conforme a los
criterios establecidos en el manual de campo
(\citeproc{ref-ministeriodedesarrollosocialyfamilia_informe_2024}{Ministerio
de Desarrollo Social y Familia, 2024, pp. 14, 56}). La muestra lograda
alcanzó 218.367 personas, pertenecientes a 78.654 hogares, con una tasa
de respuesta de 74,1\%
(\citeproc{ref-ministeriodedesarrollosocialyfamilia_informe_2024}{Ministerio
de Desarrollo Social y Familia, 2024}).

Para esta investigación, la unidad de análisis son las personas. El
estudio se realiza en dos etapas: primero se describen la población y la
utilización según grupo de edad, ingreso, previsión y territorio
mediante porcentajes ponderados, y luego se estiman modelos
multivariados. Ambas etapas se restringen a la población de 18 años y
más (\emph{N} = 172.309), porque la utilización de los menores de 18
años depende principalmente de las decisiones y los recursos de las
personas adultas del hogar.

Todas las estimaciones estadísticas se realizarán aplicando los factores
de expansión poblacional (expr) y las variables de estratos (varstrat) y
conglomerados (varunit), considerando el diseño muestral complejo de la
CASEN, para garantizar que los resultados sean representativos a nivel
nacional. Las restricciones de la muestra se realizarán mediante
estimación de subpoblaciones, y no eliminando casos, para estimar
correctamente la varianza. Las respuestas ``no sabe'' y ``no responde''
se tratarán como datos perdidos. Para que los modelos anidados sean
comparables, todos se estimarán sobre la misma muestra analítica: las
personas con información completa en todas las variables del modelo más
completo.

\section{Variables}\label{variables}

\subsection{Variables dependientes}\label{variables-dependientes}

Para esta investigación, las variables dependientes corresponden a tres
indicadores de utilización declarada de servicios de salud, construidos
como variables dicotómicas (1 = utilizó el servicio en el período de
referencia; 0 = no lo utilizó): (a) atención de medicina general en los
últimos tres meses (s20a\_preg); (b) atención de especialidad en los
últimos tres meses (s23a\_preg); y (c) hospitalización o intervención
quirúrgica en los últimos doce meses (s27a). Estos indicadores miden
utilización efectivamente declarada, y no permiten evaluar la
oportunidad, la pertinencia ni la calidad de la atención. Dado que los
períodos de referencia difieren entre servicios, la comparación entre
ellos se realiza sobre patrones y no sobre niveles de utilización.

\subsection{Variables independientes}\label{variables-independientes}

Las variables independientes se organizarán de acuerdo con las
dimensiones del Modelo Conductual de Utilización de los Servicios de
Salud de Andersen (\citeproc{ref-andersen_revisiting_1995}{1995},
\citeproc{ref-andersen_national_2008}{2008}), distinguiendo entre
factores predisponentes, facilitadores y de necesidad.

Los factores predisponentes corresponden a características
sociodemográficas que anteceden a la utilización y pueden influir en la
propensión a buscar atención. Se considerarán la edad, el sexo y el
nivel educacional (educc). La edad se agrupa en cinco tramos 18-34,
35-49, 50-64 y 65-79 y 80+. Aunque la Ley N° 19.828 define como adulta
mayor a la persona de 60 años y más, se utiliza el umbral de 65 años
para mantener la comparabilidad con las cifras del INE citadas en la
introducción, y el de 80 años corresponde a la cuarta edad que define la
misma ley. El sexo se operacionaliza como una variable categórica
nominal dicotómica, codificada como 1 = hombre y 2 = mujer. Siguiendo a
Andersen (\citeproc{ref-andersen_revisiting_1995}{1995},
\citeproc{ref-andersen_national_2008}{2008}), la educación se trata como
una característica de la estructura social y no como un recurso
material. Se operacionaliza como categórica ordinal y se agrupa en
cuatro categorías: básica incompleta o menos (códigos 0 y 1), básica
completa o media incompleta (códigos 2 y 3), media completa (código 4) y
educación superior (códigos 5 y 6). Los valores negativos, incluido el
código −88 correspondiente a ``no sabe'' y se consideran como valores
perdidos.

Los factores de necesidad corresponden a condiciones de salud que pueden
generar requerimientos de atención. Se consideran dos indicadores de
necesidad percibida, y no evaluada por criterios profesionales
(\citeproc{ref-andersen_revisiting_1995}{1995},
\citeproc{ref-andersen_national_2008}{2008}): haber tenido un problema
de salud, enfermedad o accidente en los últimos tres meses (s16), cuyo
período coincide con el de la atención de medicina general y de
especialidad, y la limitación funcional, construida a partir de seis
dominios (h7a a h7f): se considera que hay limitación si la persona
declara mucha dificultad o imposibilidad en al menos uno de ellos, según
el criterio Washington Group dado en la misma encuesta. Hay que
considerar que ambos indicadores son declarados por el informante del
hogar y no equivalen a una evaluación clínica.

Los factores facilitadores individuales comprenden los recursos
materiales y financieros que pueden favorecer o dificultar la
utilización de los servicios. Se incluirán el sistema previsional de
salud (s13), agrupado en Fonasa e Isapre, y el quintil de ingreso total
per cápita del hogar (ypchtotcor). Las personas con otro sistema
previsional (FF.AA. y del Orden, ninguno u otro) quedan fuera del
análisis, y los tramos de Fonasa (A--D) no se desagregan porque se
definen por ingreso, que ya se incluye en el modelo. Se utiliza el
ingreso per cápita porque permite comparar hogares de distinto tamaño y
porque el gasto en salud se resuelve en gran medida en el hogar, y no
solo con ingresos individuales, que no reflejan los recursos de quienes
no perciben ingresos propios. Además,se emplea el ingreso total, que
incluye las transferencias monetarias del Estado, y no el autónomo, por
ser relevantes para la población mayor. Se agrupa en quintiles,
calculados con el factor de expansión sobre la población total, porque
describen la posición relativa, permiten capturar relaciones no lineales
y reducen el efecto de los valores extremos.

Los factores facilitadores territoriales incorporan las barreras
espaciales y de disponibilidad de recursos. Se incluirán la región de
residencia como una variable categórica nominal que identifica las
distintas regiones del país el área geográfica (urbano/rural)
operacionalizada como una variable categórica nominal dicotómica,
codificada como 1 = urbano y 2 = rural, y la proximidad declarada de la
vivienda a un centro de salud (v35c) que identifica si la vivienda se
encuentra a menos de 20 cuadras o 2,5 km de un centro de salud de
atención primaria o de nivel superior, con valores 1=sí y 2= no. Hay que
considerar que esta variable constituye una aproximación a la cercanía
territorial de los servicios y no una medida directa del acceso
efectivo.

\subsection{Valores perdidos}\label{valores-perdidos}

Los porcentajes de valores perdidos se calcularon sobre los 172.309
adultos de la muestra, sin ponderar. Son de 0,0 \% en el área de
residencia; 0,1 \% en el quintil de ingreso y en la limitación
funcional; 0,2 \% en la cercanía a un centro de salud, en la atención de
medicina general y en la de especialidad; 0,3 \% en el nivel
educacional; y 0,7 \% en el problema de salud reciente. La edad, el sexo
y la región no presentan valores perdidos.

La previsión y la hospitalización tienen porcentajes mayores. En la
previsión, el 4,4 \% de los adultos queda fuera del análisis: el 4,0 \%
tiene otra previsión, que se excluye por decisión metodológica, y el 0,4
\% no tiene dato. En la hospitalización, el 1,0 \% queda fuera: el 0,8
\% corresponde a hospitalizaciones por embarazo, parto o cesárea, que se
excluyen por corresponder a una lógica distinta a la de necesidad
asociada a la edad, y el 0,2 \% no tiene dato.

\section{Estrategia de análisis}\label{estrategia-de-anuxe1lisis}

Se utilizarán modelos de regresión logística binaria para estimar la
probabilidad de utilización de cada uno de los tres tipos de servicios
de salud (medicina general, atención de especialidad y hospitalización o
intervención quirúrgica) en la población de 18 años y más. Todas las
estimaciones incorporan los factores de expansión y el diseño muestral
de CASEN 2024. Se estimará un conjunto de modelos anidados y
secuenciales para cada resultado:

\begin{itemize}
\item
  \textbf{Modelo 1: factores predisponentes.} Incorpora edad, sexo y
  nivel educacional.
\item
  \textbf{Modelo 2: factores de necesidad.} Agrega el problema de salud
  reciente y la limitación funcional.
\item
  \textbf{Modelo 3: factores facilitadores.} Agrega el ingreso y la
  previsión de salud (H1 y H2; objetivo 2) y el área de residencia y la
  cercanía a un centro de salud (H3; objetivo 3). La región de
  residencia se incluye como variable de control. Con los resultados de
  este modelo para medicina general y hospitalización se exploran además
  P1 y P2 (objetivos 2 y 3).
\item
  \textbf{Modelo 4: interacciones.} Para cada servicio, agrega por
  separado las interacciones entre el grupo de edad y el ingreso (P3),
  el territorio (área de residencia y cercanía a un centro de salud)
  (P4) y la previsión de salud (P5) (objetivo 4), y entre el sexo y
  estos factores (P6) (objetivo 5). Cada interacción se estima en un
  modelo separado.
\end{itemize}

Previo a la estimación, se evaluará la multicolinealidad entre nivel
educacional, ingreso y previsión de salud mediante el factor de
inflación de la varianza (VIF).

\bookmarksetup{startatroot}

\chapter{Resultados descriptivos}\label{resultados-descriptivos}

\textbf{Tabla 1. Caracterización de la población de 18 años y más (\%)}

La Tabla 1 presenta las características de la población de 18 años y más
según los factores predisponentes, facilitadores y de necesidad. En
términos de la estructura etaria, el 31,7\% de la población adulta tiene
entre 18 y 34 años y el 27,8\% entre 35 y 49 años, mientras que las
personas de 65 años y más representan el 17,9\%, de las cuales el 4,1\%
tiene 80 años o más. En la distribución por sexo, los hombres
representan un 48,9\% y las mujeres un 51,1\%. El nivel educacional se
distribuye principalmente en educación superior con un 41,7\%, seguida
de un 31\% con educación media completa frente a valores inferiores de
educación básica completa o media incompleta con 17,4\% y básica
incompleta o menos con un 9,9\%.

Respecto de la necesidad, el 19,1\% declara haber tenido un problema de
salud en los últimos tres meses y el 11,5\% presenta limitación
funcional. En cuanto a los factores facilitadores, Fonasa concentra la
mayoría de los casos con un 86,1\% frente a un 13,9\% en Isapres. En
cuanto al territorio, el 88,5\% reside en zonas urbanas y el 81,8\%
declara contar con un centro de salud cercano a su vivienda. Por último,
la distribución por quintiles de ingreso varía desde un 16,3\% en el
quintil de menores ingresos (Q1) hasta un 22\% en los quintiles Q4 y Q5.
Esto se debe a que los quintiles se calcularon sobre la población total
y a que los menores de 18 años se concentran en los quintiles de menores
ingresos, en parte porque el ingreso per cápita es menor en los hogares
con más integrantes.

\textbackslash begin\{table\}{[}!h{]} \centering
\textbackslash caption\{\label{tab:unnamed-chunk-2}Tabla 1.
Caracterización de la población de 18 años y más (\%)\} \centering
\resizebox{\ifdim\width>\linewidth\linewidth\else\width\fi}{!}{
\fontsize{10}{12}\selectfont
\begin{tabular}[t]{lll}
\toprule
\textbf{Variable} & \textbf{Categoría} & \textbf{\%}\\
\midrule
Grupo de edad & 18-34 & 31,7\\
 & 35-49 & 27,8\\
 & 50-64 & 22,6\\
 & 65-79 & 13,8\\
 & 80+ & 4,1\\
Sexo & Hombre & 48,9\\
 & Mujer & 51,1\\
Nivel educacional & Básica incompleta o menos & 9,9\\
 & Básica completa o media incompleta & 17,4\\
 & Media completa & 31,0\\
 & Superior & 41,7\\
Problema de salud reciente & Sí & 19,1\\
 & No & 80,9\\
Limitación funcional & Sí & 11,5\\
 & No & 88,5\\
Previsión de salud & FONASA & 86,1\\
 & Isapre & 13,9\\
Quintil de ingreso & Q1 & 16,3\\
 & Q2 & 18,8\\
 & Q3 & 20,9\\
 & Q4 & 22,0\\
 & Q5 & 22,0\\
Área de residencia & Urbano & 88,5\\
 & Rural & 11,5\\
Centro de salud a menos de 20 cuadras o 2,5 km & Sí & 81,8\\
 & No & 18,2\\
\bottomrule
\end{tabular}} \textbackslash end\{table\}

\emph{Nota.} Elaboración propia con datos de CASEN 2024.

\textbf{Tabla 2. Necesidad de salud y utilización de servicios según
grupo de edad (\%)}

La Tabla 2 muestra la distribución de los dos indicadores de necesidad y
los tres tipos de utilización de los servicios de salud según los grupos
etarios. Ambos indicadores de necesidad presentan un incremento
progresivo a medida que aumenta la edad. El aumento más notorio se
evidencia en las limitaciones funcionales, cuya prevalencia aumenta de
3,8\% en el grupo 18-34 a 54,2\% en las personas de 80 años y más. La
utilización de servicios también aumenta con la edad: en este último
grupo, el 44,7\% utilizó medicina general, el 31,9\% especialidad y el
14,7\% se hospitalizó. Con todo, el aumento de la utilización es más
moderado que el de la limitación funcional: la medicina general más que
se duplica (19,0\% a 44,7\%), la especialidad se duplica aproximadamente
(16,1\% a 31,9\%) y la hospitalización se triplica (4,7\% a 14,7\%),
mientras que la limitación funcional se multiplica por más de catorce.
Además, en las edades más altas el aumento de la utilización se hace más
lento: entre los 65-79 años y los 80 años y más la limitación funcional
se duplica (26,1\% a 54,2\%), pero la medicina general sube de 40,3\% a
44,7\% y la especialidad de 29,0\% a 31,9\%.

\textbackslash begin\{table\}{[}!h{]} \centering
\textbackslash caption\{\label{tab:unnamed-chunk-3}Tabla 2. Necesidad de
salud y utilización de servicios según grupo de edad (\%)\} \centering
\resizebox{\ifdim\width>\linewidth\linewidth\else\width\fi}{!}{
\fontsize{10}{12}\selectfont
\begin{tabular}[t]{llllll}
\toprule
\textbf{Grupo de edad} & \textbf{Problema reciente} & \textbf{Limitación funcional} & \textbf{Medicina general} & \textbf{Especialidad} & \textbf{Hospitalización}\\
\midrule
18-34 & 13,7 & 3,8 & 19,0 & 16,1 & 4,7\\
35-49 & 17,4 & 5,1 & 24,4 & 19,6 & 6,9\\
50-64 & 21,5 & 13,4 & 31,3 & 22,2 & 8,7\\
65-79 & 26,9 & 26,1 & 40,3 & 29,0 & 12,2\\
80+ & 31,7 & 54,2 & 44,7 & 31,9 & 14,7\\
Total 18 años y más & 19,1 & 11,5 & 27,3 & 20,9 & 7,7\\
\bottomrule
\end{tabular}} \textbackslash end\{table\}

\emph{Nota.} Elaboración propia con datos de CASEN 2024.

\textbf{Tabla 3. Utilización de servicios según ingreso, previsión y
territorio (\%)}

La Tabla 3 muestra el porcentaje de personas que utilizan cada servicio
según su quintil de ingreso, su previsión de salud y su territorio de
residencia. En los tres servicios la utilización aumenta con el quintil
de ingreso, y el gradiente más marcado se observa en la atención de
especialidad, que va de 14,3\% en el Q1 a 30,4\% en el Q5; en medicina
general va de 22,0\% a 31,2\% y en hospitalización de 6,5\% a 8,9\%. Las
personas afiliadas a Isapre utilizan más los tres servicios que las
afiliadas a Fonasa, con la mayor diferencia en la atención de
especialidad (33,7\% frente a 19,1\%). En cuanto al territorio, residir
en zona rural se asocia con menor utilización en los tres servicios,
especialmente en especialidad (15,8\% frente a 21,5\% en zona urbana).
La cercanía a un centro de salud muestra diferencias menores: quienes no
cuentan con uno cercano utilizan menos la medicina general (26,0\%
frente a 27,6\%) y la especialidad (18,5\% frente a 21,4\%), sin
diferencias en hospitalización (7,7\% en ambos casos). Se trata de
porcentajes brutos, que no consideran las necesidades de salud ni la
relación entre los propios factores, por lo que serán evaluados en los
modelos multivariados.

\textbackslash begin\{table\}{[}!h{]} \centering
\textbackslash caption\{\label{tab:unnamed-chunk-4}Tabla 3. Utilización
de servicios según ingreso, previsión y territorio (\%)\} \centering
\resizebox{\ifdim\width>\linewidth\linewidth\else\width\fi}{!}{
\fontsize{10}{12}\selectfont
\begin{tabular}[t]{lllll}
\toprule
\textbf{Variable} & \textbf{Categoría} & \textbf{Medicina general} & \textbf{Especialidad} & \textbf{Hospitalización}\\
\midrule
Quintil de ingreso & Q1 & 22,0 & 14,3 & 6,5\\
 & Q2 & 25,4 & 16,8 & 7,1\\
 & Q3 & 27,4 & 18,6 & 7,7\\
 & Q4 & 28,9 & 21,9 & 7,8\\
 & Q5 & 31,2 & 30,4 & 8,9\\
Previsión de salud & FONASA & 27,1 & 19,1 & 7,4\\
 & Isapre & 30,6 & 33,7 & 9,7\\
Área de residencia & Urbano & 27,6 & 21,5 & 7,8\\
 & Rural & 25,0 & 15,8 & 6,8\\
Centro de salud cercano & Sí & 27,6 & 21,4 & 7,7\\
 & No & 26,0 & 18,5 & 7,7\\
\bottomrule
\end{tabular}} \textbackslash end\{table\}

\emph{Nota.} Elaboración propia con datos de CASEN 2024.

\bookmarksetup{startatroot}

\chapter{Referencias
bibliográficas}\label{referencias-bibliogruxe1ficas}

\protect\phantomsection\label{refs}
\begin{CSLReferences}{1}{0}
\bibitem[\citeproctext]{ref-aday_framework_1974}
Aday, L. A., \& Andersen, R. (1974). A framework for the study of access
to medical care. \emph{Health Services Research}, \emph{9}(3), 208-220.
\url{https://pmc.ncbi.nlm.nih.gov/articles/instance/1071804/pdf/hsresearch00560-0030.pdf}

\bibitem[\citeproctext]{ref-albala_envejecimiento_2020}
Albala, C. (2020). El envejecimiento de la población chilena y los
desafíos para la salud y el bienestar de las personas mayores.
\emph{Revista Médica Clínica Las Condes}, \emph{31}(1), 7-12.
\url{https://doi.org/10.1016/j.rmclc.2019.12.001}

\bibitem[\citeproctext]{ref-andersen_revisiting_1995}
Andersen, R. M. (1995). Revisiting the behavioral model and access to
medical care: {Does} it matter? \emph{Journal of Health and Social
Behavior}, \emph{36}(1), 1-10. \url{https://doi.org/10.2307/2137284}

\bibitem[\citeproctext]{ref-andersen_national_2008}
Andersen, R. M. (2008). National health surveys and the behavioral model
of health services use. \emph{Medical Care}, \emph{46}(7), 647-653.
\url{https://doi.org/10.1097/MLR.0b013e31817a835d}

\bibitem[\citeproctext]{ref-asadi-lari_need_2003}
Asadi-Lari, M., Packham, C., \& Gray, D. (2003). Need for redefining
needs. \emph{Health and Quality of Life Outcomes}, \emph{1}, Artículo
34. \url{https://doi.org/10.1186/1477-7525-1-34}

\bibitem[\citeproctext]{ref-bancomundial_camas_}
Banco Mundial. (s.~f.-a). \emph{Camas hospitalarias (cada 1.000
personas) - Chile}. Recuperado 18 de octubre de 2026, de
\url{https://datos.bancomundial.org/indicator/SH.MED.BEDS.ZS?locations=CL}

\bibitem[\citeproctext]{ref-bancomundial_medicos_}
Banco Mundial. (s.~f.-b). \emph{Médicos (cada 1.000 personas) - Chile}.
Recuperado 18 de octubre de 2026, de
\url{https://datos.bancomundial.org/indicador/SH.MED.PHYS.ZS?locations=CL}

\bibitem[\citeproctext]{ref-bassdelcampo_universalizacion_2025}
Bass del Campo, C. (2025). Universalización de Atención Primaria en
Chile: una política lamentablemente regresiva. \emph{Cuadernos Médico
Sociales}, \emph{65}(1), 53-65.
\url{https://doi.org/10.56116/cms.v65.n1.2025.2209}

\bibitem[\citeproctext]{ref-bibliotecadelcongresonacionaldechile_ley_}
Biblioteca del Congreso Nacional de Chile. (s.~f.). \emph{Ley de
urgencia en salud} {[}Ley Fácil{]}. Recuperado 27 de octubre de 2026, de
\url{https://www.bcn.cl/portal/leyfacil/recurso/ley-de-urgencia-en-salud}

\bibitem[\citeproctext]{ref-blanco_enfoque_2011}
Blanco, M. (2011). El enfoque del curso de vida: orígenes y desarrollo.
\emph{Revista Latinoamericana de Población}, \emph{5}(8), 5-31.
\url{https://doi.org/10.31406/relap2011.v5.i1.n8.1}

\bibitem[\citeproctext]{ref-chackiel_envejecimiento_2000}
Chackiel, J. (2000). \emph{El envejecimiento de la población
latinoamericana: ¿hacia una relación de dependencia favorable?} Comisión
Económica para América Latina y el Caribe.
\url{https://repositorio.cepal.org/server/api/core/bitstreams/c15fd12a-15b2-4b61-91f6-4863270ae3d8/content}

\bibitem[\citeproctext]{ref-chovarvera_desigualdad_2014}
Chovar Vera, A., Vásquez Lavín, F., \& Paraje, G. (2014). Desigualdad e
inequidad en la utilización de servicios médicos según grupos etarios en
Chile, 2000--2011. \emph{Revista Panamericana de Salud Pública},
\emph{36}(3), 171-178.
\url{https://www.scielosp.org/pdf/rpsp/2014.v36n3/171-178/es}

\bibitem[\citeproctext]{ref-clinicasdechile_dimensionamiento_2025}
Clínicas de Chile. (2025). \emph{Dimensionamiento de la oferta de salud
en Chile 2024}. Asociación de Clínicas de Chile AG.
\url{https://www.clinicasdechile.cl/wp-content/uploads/2025/12/Dimensionamiento-del-Sector-Salud-en-Chile-\%E2\%80\%93-Cifras-a-2024.pdf}

\bibitem[\citeproctext]{ref-comisioneconomicaparaamericalatinayelcaribe_acerca_}
Comisión Económica para América Latina y el Caribe. (s.~f.).
\emph{Acerca de envejecimiento}. Recuperado
\url{https://www.cepal.org/es/temas/envejecimiento/acerca-envejecimiento}

\bibitem[\citeproctext]{ref-comisioneconomicaparaamericalatinayelcaribe_impactos_2026}
Comisión Económica para América Latina y el Caribe. (2026).
\emph{Impactos del cambio demográfico en América Latina y el Caribe:
Retos y opciones para la política pública} (LC/CRPD.6/3).

\bibitem[\citeproctext]{ref-dannefer_cumulative_2003}
Dannefer, D. (2003). Cumulative advantage/disadvantage and the life
course: {Cross}-fertilizing age and social science theory. \emph{The
Journals of Gerontology Series B: Psychological Sciences and Social
Sciences}, \emph{58}(6), S327-S337.
\url{https://doi.org/10.1093/geronb/58.6.S327}

\bibitem[\citeproctext]{ref-departamentodeestadisticaseinformaciondesalud_defunciones_2026}
Departamento de Estadísticas e Información de Salud. (2026, mayo 5).
\emph{Defunciones y mortalidad por capítulo CIE-10 de causa básica de
defunción y año} {[}Tablero de datos{]}. Ministerio de Salud de Chile.
\url{https://informesdeis.minsal.cl/SASVisualAnalytics/?reportUri=\%2Freports\%2Freports\%2F4013de47-a3c2-47b8-8547-075525e4f819&sectionIndex=0&sso_guest=true&reportViewOnly=true&reportContextBar=false&sas-welcome=false}

\bibitem[\citeproctext]{ref-equipoeunamed_como_2025}
Equipo EUNAMed. (2025, octubre 10). \emph{Cómo funciona el sistema de
salud en Chile}. EUNAMed.com.
\url{https://eunamed.com/blog/como-funciona-el-sistema-de-salud-en-chile/}

\bibitem[\citeproctext]{ref-ferraro_aging_2009}
Ferraro, K. F., \& Shippee, T. P. (2009). Aging and cumulative
inequality: {How} does inequality get under the skin? \emph{The
Gerontologist}, \emph{49}(3), 333-343.
\url{https://doi.org/10.1093/geront/gnp034}

\bibitem[\citeproctext]{ref-fuenzalida_analisis_2023}
Fuenzalida, M., Crespo, F., \& Lanio, I. (2023). Análisis de la
cobertura geográfica de camas hospitalarias públicas en Chile.
\emph{Revista Geográfica de Chile Terra Australis}, \emph{59}, 92-101.
\url{https://doi.org/10.23854/07199562.2023591.fuenzalida}

\bibitem[\citeproctext]{ref-gattini_sistema_2018}
Gattini, C. (2018). \emph{El sistema de salud en Chile}. Observatorio
Chileno de Salud Pública, Universidad de Chile.
\url{https://www.ochisap.cl/wp-content/uploads/2022/04/Sistema_Salud_Chile_Gattini_2018.pdf}

\bibitem[\citeproctext]{ref-goldstein_sistema_2018}
Goldstein, E. (2018). \emph{El sistema de salud en Chile y la Atención
Primaria de Salud municipal: Marco para un debate sobre
desmunicipalización}. Biblioteca del Congreso Nacional de Chile.
\url{https://obtienearchivo.bcn.cl/obtienearchivo?id=repositorio/10221/26811/2/BCN_Gobernanza_salud_y_demunicip_para_reposit_final.pdf}

\bibitem[\citeproctext]{ref-goldstein_resena_2024}
Goldstein, E. (2024). \emph{Reseña sobre la Ley Corta de ISAPRES de 2024
(Ley N° 21.674)}. Biblioteca del Congreso Nacional de Chile.
\url{https://obtienearchivo.bcn.cl/obtienearchivo?id=repositorio/10221/36273/1/BCN__Resena_Ley_Corta_Isapres_2024__final.pdf}

\bibitem[\citeproctext]{ref-gonzalez_estructura_2019}
González, C., Castillo-Laborde, C., \& Matute, I. (Eds.). (2019).
\emph{Estructura y funcionamiento del sistema de salud chileno} (Serie
de Salud Poblacional N.º 2). Universidad del Desarrollo.
\url{http://hdl.handle.net/11447/2895}

\bibitem[\citeproctext]{ref-henriquez_desigualdades_2015}
Henríquez, J., \& Velasco, C. (2015). \emph{Las desigualdades en la
atención médica en los últimos 20 años} (Propuestas de Política Pública
N.° 13). Centro de Estudios Públicos.

\bibitem[\citeproctext]{ref-huaiquilaf-jorquera_medicos_2026}
Huaiquilaf-Jorquera, S., Riedemann Gonzalez, J. P., Castillo Mora, J.
A., \& Arancibia Luna, M. J. (2026). Médicos especialistas en Chile:
Análisis de la situación 2024. \emph{Revista Médica de Chile},
\emph{154}(3).
\url{https://www.revistamedicadechile.cl/index.php/rmedica/article/view/11531}

\bibitem[\citeproctext]{ref-huenchuan_proteccion_2011}
Huenchuan, S. (2011). \emph{La protección de la salud en el marco de la
dinámica demográfica y los derechos} (Serie Población y Desarrollo).
Comisión Económica para América Latina y el Caribe.

\bibitem[\citeproctext]{ref-huenchuan_envejecimiento_2016}
Huenchuan, S. (Ed.). (2016). \emph{Envejecimiento e institucionalidad
pública en América Latina y el Caribe: conceptos, metodologías y casos
prácticos}. Comisión Económica para América Latina y el Caribe.
\url{https://www.cepal.org/es/publicaciones/40197-envejecimiento-institucionalidad-publica-america-latina-caribe-conceptos}

\bibitem[\citeproctext]{ref-huenchuan_envejecimiento_2018}
Huenchuan, S. (Ed.). (2018). \emph{Envejecimiento, personas mayores y
Agenda 2030 para el Desarrollo Sostenible: perspectiva regional y de
derechos humanos}. Comisión Económica para América Latina y el Caribe.
\url{https://repositorio.cepal.org/bitstream/handle/11362/44369/1/S1800629_es.pdf}

\bibitem[\citeproctext]{ref-rojas_envejecimiento_}
Hugo Rojas, F., Rodríguez Canache, L., \& Rodríguez León, J. (2022).
\emph{Envejecimiento en Chile: Evolución, características de las
personas mayores y desafíos demográficos para la población} (Documentos
de Trabajo). Instituto Nacional de Estadísticas.
\url{https://www.ine.gob.cl/docs/default-source/documentos-de-trabajo/envejecimiento-en-chile-evolucion-y-caracteristicas-de-las-personas-mayores.pdf?sfvrsn=fa394551_2}

\bibitem[\citeproctext]{ref-institutonacionaldeestadisticas_transicion_2022}
Instituto Nacional de Estadísticas. (2022). \emph{Transición demográfica
y el envejecimiento de la población en Chile: Evolución, características
de las personas mayores y desafíos demográficos para la población}
{[}Presentación{]}. Conferencia ciudadana, Subdepartamento de
Demografía.
\url{https://www.ine.gob.cl/docs/default-source/conferencias/evolucion-demografica/conferencias-ciudadanas-evolucion-demografica-en-chile.pdf}

\bibitem[\citeproctext]{ref-institutonacionaldeestadisticas_cifras_2025}
Instituto Nacional de Estadísticas. (2025a). \emph{Cifras provisionales
de estadísticas vitales confirman disminución de nacimientos en 2024}
{[}Noticia{]}.
\url{https://www.ine.gob.cl/sala-de-prensa/prensa/general/noticia/2025/05/15/cifras-provisionales-de-estad\%C3\%ADsticas-vitales-confirman-disminuci\%C3\%B3n-de-nacimientos-en-2024}

\bibitem[\citeproctext]{ref-institutonacionaldeestadisticas_primeros_2025}
Instituto Nacional de Estadísticas. (2025b). \emph{Primeros resultados
del Censo 2024: 18.480.432 personas fueron censadas en Chile,
manteniéndose la tendencia de envejecimiento de la población}.
\url{https://censo2024.ine.gob.cl/primeros-resultados-del-censo-2024-18-480-432-personas-fueron-censadas-en-chile-manteniendose-la-tendencia-de-envejecimiento-de-la-poblacion/}

\bibitem[\citeproctext]{ref-institutonacionaldeestadisticas_estimaciones_2026}
Instituto Nacional de Estadísticas. (2026). \emph{Estimaciones y
proyecciones de población 1992-2070: Base 2024. Presentación de
resultados}.
\url{https://www.ine.gob.cl/docs/default-source/proyecciones-de-poblacion/publicaciones-y-anuarios/proyecci\%C3\%B3n-base-2024/estimaciones-y-proyecciones-de-poblaci\%C3\%B3n-1992-2070_base-2024_presentaci\%C3\%B3n-de-resultados.pdf}

\bibitem[\citeproctext]{ref-lantz_socioeconomic_2001}
Lantz, P. M. (2001). Socioeconomic status and health care. En N. J.
Smelser \& P. B. Baltes (Eds.), \emph{International encyclopedia of the
social \& behavioral sciences} (pp. 14558-14562). Elsevier.
\url{https://doi.org/10.1016/B0-08-043076-7/03828-6}

\bibitem[\citeproctext]{ref-levesque_patientcentred_2013}
Levesque, J.-F., Harris, M. F., \& Russell, G. (2013). Patient-centred
access to health care: {Conceptualising} access at the interface of
health systems and populations. \emph{International Journal for Equity
in Health}, \emph{12}, Artículo 18.
\url{https://doi.org/10.1186/1475-9276-12-18}

\bibitem[\citeproctext]{ref-_ley_2002}
Ley N° 19.828. (2002, septiembre 27). \emph{Crea el Servicio Nacional
del Adulto Mayor}. Ministerio Secretaría General de la Presidencia.
\url{https://www.bcn.cl/leychile/navegar?idNorma=202950}

\bibitem[\citeproctext]{ref-mendoza-sassi_utilizacion_2001}
Mendoza-Sassi, R., \& Béria, J. U. (2001). Utilización de los servicios
de salud: una revisión sistemática sobre los factores relacionados.
\emph{Cadernos de Saúde Pública}, \emph{17}(4), 819-832.
\url{https://www.scielosp.org/article/csp/2001.v17n4/819-832/}

\bibitem[\citeproctext]{ref-ministeriodedesarrollosocialyfamilia_informe_2024}
Ministerio de Desarrollo Social y Familia. (2024). \emph{Informe
metodológico Casen 2024}. Observatorio Social.
\url{https://observatorio.ministeriodesarrollosocial.gob.cl/storage/docs/casen/2024/Informe_metodologico_casen_2024.pdf}

\bibitem[\citeproctext]{ref-ministeriodesalud_funciones_}
Ministerio de Salud. (s.~f.). \emph{Funciones y objetivos}. Gobierno de
Chile. Recuperado 20 de octubre de 2026, de
\url{https://www.minsal.cl/funciones-objetivos/}

\bibitem[\citeproctext]{ref-ministeriodesaluddepartamentodeepidemiologia_informe_2023}
Ministerio de Salud, Departamento de Epidemiología. (2023).
\emph{Informe vigilancia enfermedades no transmisibles 2023}.
\url{https://epi.minsal.cl/wp-content/uploads/2025/05/INFORME_NACIONAL_VENT_CHILE.pdf}

\bibitem[\citeproctext]{ref-ministeriodesaluddepartamentodeepidemiologia_encuesta_2025}
Ministerio de Salud, Departamento de Epidemiología. (2025).
\emph{Encuesta Nacional de Calidad de Vida y Salud (ENCAVI) 2023-2024:
Resultados} {[}Presentación de diapositivas{]}.
\url{https://epi.minsal.cl/wp-content/uploads/2026/07/PPT_ENCAVI_2023_24.pdf}

\bibitem[\citeproctext]{ref-ministeriodesalud_levantamiento_2025}
Ministerio de Salud, \& Organización Panamericana de la Salud. (2025).
\emph{Levantamiento de información de Postas de Salud Rural}.
\url{https://www.paho.org/sites/default/files/2026/09/minsal-ops-2025-levantamiento-informacion-psr-final-comprimido.pdf}

\bibitem[\citeproctext]{ref-moscosocuaresma_determinantes_2017}
Moscoso Cuaresma, J. R., Cárdenas, E., \& Juárez, C. (2017).
\emph{Determinantes sociales en salud}. Universidad Peruana de Ciencias
Aplicadas.
\url{https://www.researchgate.net/profile/Julio-Ricardo-Moscoso-Cuaresma/publication/385822723_Determinantes_sociales_en_salud/links/6736a966a78ba469f0617e2a/Determinantes-sociales-en-salud.pdf}

\bibitem[\citeproctext]{ref-nunez_equity_2013}
Núñez, A., \& Chi, C. (2013). Equity in health care utilization in
{Chile}. \emph{International Journal for Equity in Health}, \emph{12},
Artículo 58. \url{https://doi.org/10.1186/1475-9276-12-58}

\bibitem[\citeproctext]{ref-observatoriodelenvejecimientouc-confuturo_envejecer_2023}
Observatorio del Envejecimiento UC-Confuturo. (2023). \emph{Envejecer en
la ruralidad: Explorando una realidad invisibilizada} (Año 5, vol. 29).

\bibitem[\citeproctext]{ref-olguin_niveles_2025}
Olguín, C. (2025, octubre 20). \emph{Niveles de atención en salud: El
rol de la atención secundaria y terciaria}. ADIPA.
\url{https://adipa.cl/noticias/niveles-de-atencion-en-salud-el-rol-de-la-atencion-secundaria-y-terciaria/}

\bibitem[\citeproctext]{ref-organisationforeconomicco-operationanddevelopment_health_2025}
Organisation for Economic Co-operation and Development. (2025).
\emph{{Health at a Glance 2025}: {Country} notes---{Chile}}.
\url{https://www.oecd.org/es/publications/2025/11/health-at-a-glance-2025-country-notes_2f94481e/chile_b440e042.html}

\bibitem[\citeproctext]{ref-paraje_health_2012}
Paraje, G., \& Vásquez, F. (2012). Health equity in an unequal country:
{The} use of medical services in {Chile}. \emph{International Journal
for Equity in Health}, \emph{11}, Artículo 81.
\url{https://doi.org/10.1186/1475-9276-11-81}

\bibitem[\citeproctext]{ref-penchansky_concept_1981}
Penchansky, R., \& Thomas, J. W. (1981). The concept of access:
{Definition} and relationship to consumer satisfaction. \emph{Medical
Care}, \emph{19}(2), 127-140.
\url{https://doi.org/10.1097/00005650-198102000-00001}

\bibitem[\citeproctext]{ref-programadelasnacionesunidasparaeldesarrollo_desigualdad_2018}
Programa de las Naciones Unidas para el Desarrollo. (2018).
\emph{Desigualdad regional en Chile: Ingresos, salud y educación en
perspectiva territorial}.

\bibitem[\citeproctext]{ref-proyectonodo_comunidades_2021}
Proyecto NODO. (2021). \emph{Comunidades y entornos rurales en Chile.
Análisis y perspectivas sobre población mayor}. Sistema de las Naciones
Unidas en Chile.

\bibitem[\citeproctext]{ref-robledomarin_teorias_2021}
Robledo Marín, C. A., \& Orejuela Gómez, J. J. (2021). Teorías de la
sociología del envejecimiento y la vejez. \emph{Revista Guillermo de
Ockham}, \emph{18}(1), 95-102.
\url{https://doi.org/10.21500/22563202.4660}

\bibitem[\citeproctext]{ref-rodriguezsantana_need_2023}
Rodriguez Santana, I., Mason, A., Gutacker, N., Kasteridis, P., Santos,
R., \& Rice, N. (2023). Need, demand, supply in health care: {Working}
definitions, and their implications for defining access. \emph{Health
Economics, Policy and Law}, \emph{18}(1), 1-13.
\url{https://doi.org/10.1017/S1744133121000293}

\bibitem[\citeproctext]{ref-serviciodesaludmetropolitanosur_rol_}
Servicio de Salud Metropolitano Sur. (s.~f.). \emph{Rol de la red
asistencial}. Ministerio de Salud, Gobierno de Chile. Recuperado 18 de
octubre de 2026, de \url{https://ssms.gob.cl/rol-de-la-red-asistencial/}

\bibitem[\citeproctext]{ref-superintendenciadesalud_como_}
Superintendencia de Salud. (s.~f.-a). \emph{Cómo funciona el Sistema de
Salud en Chile}. Ministerio de Salud, Gobierno de Chile. Recuperado 20
de octubre de 2026, de
\url{https://www.superdesalud.gob.cl/orientacion-en-salud/como-funciona-el-sistema-de-salud-en-chile/}

\bibitem[\citeproctext]{ref-superintendenciadesalud_garantias_}
Superintendencia de Salud. (s.~f.-b). \emph{Garantías Explícitas en
Salud (GES)}. Gobierno de Chile. Recuperado 20 de octubre de 2026, de
\url{https://www.superdesalud.gob.cl/tax-temas-de-orientacion/garantias-explicitas-en-salud-ges-1962/}

\bibitem[\citeproctext]{ref-superintendenciadesalud_caracterizacion_2023}
Superintendencia de Salud. (2023). \emph{Caracterización de las
especialidades médicas en Chile en 2022}. Ministerio de Salud.
\url{https://www.superdesalud.gob.cl/documentacion/666/articles-26547_recurso_1.pdf}

\bibitem[\citeproctext]{ref-travassos_uma_2004}
Travassos, C., \& Martins, M. (2004). Uma revisão sobre os conceitos de
acesso e utilização de serviços de saúde. \emph{Cadernos de Saúde
Pública}, \emph{20}(Supl. 2), S190-S198.
\url{https://www.scielo.br/j/csp/a/PkyrsjDrZWwzzPVJJPbbXtQ/}

\bibitem[\citeproctext]{ref-turra_transicion_2021}
Turra, C. M., \& Fernandes, F. (2021). \emph{La transición demográfica:
Oportunidades y desafíos en la senda hacia el logro de los Objetivos de
Desarrollo Sostenible en América Latina y el Caribe} (LC/TS.2020/105).
Comisión Económica para América Latina y el Caribe.
\url{https://www.cepal.org/es/publicaciones/46805-la-transicion-demografica-oportunidades-desafios-la-senda-logro-objetivos}

\bibitem[\citeproctext]{ref-unitednationsdepartmentofeconomicandsocialaffairspopulationdivision_world_2020}
United Nations. (2020). \emph{World population ageing 2019}
(ST/ESA/SER.A/444).
\url{https://digitallibrary.un.org/record/3907988/files/WorldPopulationAgeing2019-Report.pdf}

\bibitem[\citeproctext]{ref-unitednationsdepartmentofeconomicandsocialaffairspopulationdivision_world_2022}
United Nations. (2022). \emph{World population prospects 2022: {Summary}
of results} (UN DESA/POP/2022/TR/NO. 3).
\url{https://www.un.org/development/desa/pd/sites/www.un.org.development.desa.pd/files/wpp2022_summary_of_results.pdf}

\bibitem[\citeproctext]{ref-vasquez_incomerelated_2013}
Vásquez, F., Paraje, G., \& Estay, M. (2013). Income-related inequality
in health and health care utilization in {Chile}, 2000--2009.
\emph{Revista Panamericana de Salud Pública}, \emph{33}(2), 98-106.
\url{https://doi.org/10.1590/S1020-49892013000200004}

\bibitem[\citeproctext]{ref-vignolo_niveles_2011}
Vignolo, J., Vacarezza, M., Álvarez, C., \& Sosa, A. (2011). Niveles de
atención, de prevención y atención primaria de la salud. \emph{Archivos
de Medicina Interna}, \emph{33}(1), 7-11.
\url{http://www.scielo.edu.uy/scielo.php?script=sci_arttext&pid=S1688-423X2011000100003&lng=es&tlng=es}

\end{CSLReferences}


\backmatter

% --- Cuerpo del libro (capítulos) en arábigos ---
\mainmatter
\pagestyle{scrheadings}

% --- (Opcional) Estilo del índice general ---
% \addtocontents{toc}{\protect\thispagestyle{plain}}


\end{document}
